% Generated mechanically from the frozen stacks-properties.tex target.
% Do not edit this generated file; edit the bound locale source instead.

% OK, start here.
%

\setcounter{chapter}{99}
\chapter{代数叠的性质}



\phantomsection
\label{stacks-properties-section-phantom}


\section{引言}
\label{stacks-properties-section-introduction}

\noindent
关于代数叠的简要介绍，见《代数叠》第
\ref{algebraic-section-introduction} 节；也请阅读该章的部分内容，
作为我们所用代数叠理论的基础。该章有意仔细区分概形、代数空间和代数叠；
从本章开始，我们采用通常的术语滥用，把这些概念交替使用。

\medskip\noindent
本章旨在介绍代数叠的一些基本概念和性质。
对于对角态射可表示的拟分离代数叠，一个基础参考文献是 \cite{LM-B}。



\section{约定与术语滥用}
\label{stacks-properties-section-conventions}

\noindent
我们选定一个大 fppf 位点 $\Sch_{fppf}$。
所有概形都包含在 $\Sch_{fppf}$ 中。
所考虑的每个环 $A$ 都满足：$\Spec(A)$（同构于）此大位点的一个对象。

\medskip\noindent
我们还固定一个基概形 $S$；按上述约定，它是 $\Sch_{fppf}$ 的一个元素。
只关心绝对情形的读者可以取 $S = \Spec(\mathbf{Z})$。

\medskip\noindent
关于代数叠，我们采用如下约定：
\begin{enumerate}
\item 所谓{\it 代数叠}，是指 $S$ 上的代数叠，即群胚上的纤维化范畴
$p : \mathcal{X} \to (\Sch/S)_{fppf}$
，它满足《代数叠》定义 \ref{algebraic-definition-algebraic-stack}
中的条件。
\item 称 $f : \mathcal{X} \to \mathcal{Y}$ 为{\it 代数叠的态射}，
是指 $S$ 上代数叠的 $1$-态射，即 $(\Sch/S)_{fppf}$ 上
群胚纤维化范畴的 $1$-态射；见《代数叠》定义
\ref{algebraic-definition-morphism-algebraic-stacks}。
\item 所谓 {\it $2$-态射} $\alpha : f \to g$，是指
$S$ 上代数叠的 $2$-范畴中的 $2$-态射；见《代数叠》定义
\ref{algebraic-definition-morphism-algebraic-stacks}。
\item 给定代数叠的态射 $\mathcal{X} \to \mathcal{Z}$ 和
$\mathcal{Y} \to \mathcal{Z}$，我们滥用术语，把 $2$-纤维积
$\mathcal{X} \times_\mathcal{Z} \mathcal{Y}$ 称为{\it 纤维积}。
\item 用 $\mathcal{X} \times_S \mathcal{Y}$ 表示代数叠
$\mathcal{X}$ 与 $\mathcal{Y}$ 的积。
\item 我们常滥用记号：若两个代数叠 $\mathcal{X}$ 和 $\mathcal{Y}$
在此 $2$-范畴中等价，就称它们{\it 同构}。
\end{enumerate}
关于代数空间，我们采用如下约定：
\begin{enumerate}
\item 若称 $X$ 是{\it 代数空间}，是指 $X$ 是 $S$ 上的代数空间，
即 $X$ 是 $(\Sch/S)_{fppf}$ 上满足《空间》定义
\ref{spaces-definition-algebraic-space} 中条件的预层。
\item {\it 代数空间的态射} $f :X \to Y$，是指《空间》定义
\ref{spaces-definition-morphism-algebraic-spaces} 所定义的
$S$ 上代数空间的态射。
\item 我们{\bf 不}区分代数空间 $X$ 与它所给出的代数叠
$\mathcal{S}_X \to (\Sch/S)_{fppf}$；见《代数叠》引理
\ref{algebraic-lemma-representable-algebraic}。
\item 特别地，从 $X$ 到代数叠 $\mathcal{Y}$ 的{\it 态射}
$f : X \to \mathcal{Y}$，是指代数叠的态射
$f : \mathcal{S}_X \to \mathcal{Y}$。态射 $\mathcal{Y} \to X$ 亦然。
\item 此外，给定代数叠 $\mathcal{X}$，称
{\it $\mathcal{X}$ 是代数空间}，是指 $\mathcal{X}$ 可由代数空间表示；
见《代数叠》定义
\ref{algebraic-definition-representable-by-algebraic-space}。
\item 我们采用如下记号约定：若用罗马体大写字母
（如 $X, Y, Z, A, B, \ldots$）表示一个代数叠，
则其惰性叠平凡，因而它是代数空间；见《代数叠》命题
\ref{algebraic-proposition-algebraic-stack-no-automorphisms}。
\end{enumerate}
关于概形，我们采用如下约定：
\begin{enumerate}
\item 若称 $X$ 是{\it 概形}，是指 $X$ 是 $S$ 上的概形，
即 $X$ 是 $(\Sch/S)_{fppf}$ 的对象。
\item 所谓{\it 概形的态射}，是指 $S$ 上概形的态射。
\item 我们{\bf 不}区分概形 $X$ 与它所给出的代数叠
$\mathcal{S}_X \to (\Sch/S)_{fppf}$；见《代数叠》引理
\ref{algebraic-lemma-representable-algebraic}。
\item 特别地，从概形 $X$ 到代数叠 $\mathcal{Y}$ 的{\it 态射}
$f : X \to \mathcal{Y}$，是指代数叠的态射
$f : \mathcal{S}_X \to \mathcal{Y}$。态射 $\mathcal{Y} \to X$ 亦然。
\item 此外，给定代数叠 $\mathcal{X}$，称
{\it $\mathcal{X}$ 是概形}，是指 $\mathcal{X}$ 可表示；
见《代数叠》第 \ref{algebraic-section-representable} 节。
\end{enumerate}
关于代数叠的态射，我们采用如下约定：
\begin{enumerate}
\item 若对每个概形 $T$ 和态射 $T \to \mathcal{Y}$，纤维积
$T \times_\mathcal{Y} \mathcal{X}$ 都是概形，则称代数叠的态射
$f : \mathcal{X} \to \mathcal{Y}$ {\it 可表示}，或
{\it 可由概形表示}。见《代数叠》第
\ref{algebraic-section-representable-morphism} 节。
\item 若对每个概形 $T$ 和态射 $T \to \mathcal{Y}$，纤维积
$T \times_\mathcal{Y} \mathcal{X}$ 都是代数空间，则称代数叠的态射
$f : \mathcal{X} \to \mathcal{Y}$ {\it 可由代数空间表示}。
见《代数叠》定义
\ref{algebraic-definition-representable-by-algebraic-spaces}。
在此情形，只要 $Z \to \mathcal{Y}$ 是源为代数空间的态射，
$Z \times_\mathcal{Y} \mathcal{X}$ 就是代数空间；见《代数叠》引理
\ref{algebraic-lemma-base-change-by-space-representable-by-space}。
\item 我们可以滥用术语：若代数叠的一个图在代数叠的
$2$-范畴中 $2$-交换，就称该图交换。
\end{enumerate}
注意，从代数空间到代数叠的每个态射 $X \to \mathcal{Y}$
都可由代数空间表示；见《代数叠》引理
\ref{algebraic-lemma-representable-diagonal}。
以后将不加说明地使用这一基本结果。




\section{可由代数空间表示的态射的性质}
\label{stacks-properties-section-properties-morphisms}

\noindent
我们将在另一章研究代数叠（任意）态射的性质。
对于可由代数空间表示的态射，我们知道满射、光滑、\'etale 等性质的含义。
这尤其适用于从代数空间到代数叠的态射 $X \to \mathcal{Y}$。
本节回顾其机制，列出适用的性质，并证明几个简单引理。

\medskip\noindent
第一个引理说明：若一个态射经某个平坦、局部有限表示且满射的态射作基变换后
可由代数空间表示，则它本身也可由代数空间表示。

\begin{lemma}
\label{stacks-properties-lemma-check-representable-covering}
设 $f : \mathcal{X} \to \mathcal{Y}$ 是代数叠的态射。
设 $W$ 是代数空间，且 $W \to \mathcal{Y}$ 满射、局部有限表示并平坦。
下列条件等价：
\begin{enumerate}
\item $f$ 可由代数空间表示；
\item $W \times_\mathcal{Y} \mathcal{X}$ 是代数空间。
\end{enumerate}
\end{lemma}

\begin{proof}
(1) $\Rightarrow$ (2) 是《代数叠》引理
\ref{algebraic-lemma-base-change-by-space-representable-by-space}。
反之，设 $W \to \mathcal{Y}$ 如 (2) 所述。为证明 (1)，
只需证明 $f$ 在纤维范畴上忠实；见《代数叠》引理
\ref{algebraic-lemma-characterize-representable-by-algebraic-spaces}。
条件 (2) 尤其蕴含 $W \times_\mathcal{Y} \mathcal{X} \to W$ 忠实。
于是由《叠》引理 \ref{stacks-lemma-faithful-descent}，$f$ 也忠实。
\end{proof}

\noindent
设 $P$ 是代数空间态射的一个性质，它在靶上 fppf 局部，
且在任意基变换下保持。设 $f : \mathcal{X} \to \mathcal{Y}$
是可由代数空间表示的代数叠态射。我们称
{\it $f$ 具有性质 $P$}，当且仅当对每个概形 $T$ 及态射
$T \to \mathcal{Y}$，代数空间的态射
$T \times_\mathcal{Y} \mathcal{X} \to T$ 具有性质 $P$；
见《代数叠》定义 \ref{algebraic-definition-relative-representable-property}。

\medskip\noindent
事实上，若 $f : \mathcal{X} \to \mathcal{Y}$ 可由代数空间表示且具有性质
$P$，则对代数叠的任意态射 $\mathcal{Y}' \to \mathcal{Y}$，
基变换
$\mathcal{Y}' \times_\mathcal{Y} \mathcal{X} \to \mathcal{Y}'$
也具有性质 $P$；见《代数叠》引理
\ref{algebraic-lemma-base-change-representable-by-spaces} 和
\ref{algebraic-lemma-base-change-representable-transformations-property}。
若性质 $P$ 在复合下保持，则同一结论也适用于可由代数空间表示的
代数叠态射；见《代数叠》引理
\ref{algebraic-lemma-composition-representable-by-spaces} 和
\ref{algebraic-lemma-composition-representable-transformations-property}。
此外，在此情形，可由代数空间表示且具有性质 $\mathcal{P}$ 的态射的积
$\mathcal{X}_1 \times \mathcal{X}_2 \to \mathcal{Y}_1 \times \mathcal{Y}_2$
也具有性质 $\mathcal{P}$；见《代数叠》引理
\ref{algebraic-lemma-product-representable-transformations-property}.

\medskip\noindent
最后，设 $P, P'$ 是代数空间态射的两个性质，它们在靶上 fppf 局部，
且在任意基变换下保持。若对每个态射 $f$ 都有
$P(f) \Rightarrow P'(f)$，则对可由代数空间表示的代数叠态射，
相应性质之间也有同一蕴含；见《代数叠》引理
\ref{algebraic-lemma-representable-transformations-property-implication}.
下文及后续各章将{\it 不加说明地}使用这一点。

\medskip\noindent
上述讨论适用于代数空间态射的下列各个性质：
\begin{enumerate}
\item 拟紧；见《空间的态射》引理
\ref{spaces-morphisms-lemma-base-change-quasi-compact}
以及《空间上的下降》引理
\ref{spaces-descent-lemma-descending-property-quasi-compact}；
\item 拟分离；见《空间的态射》引理
\ref{spaces-morphisms-lemma-base-change-separated}
以及《空间上的下降》引理
\ref{spaces-descent-lemma-descending-property-quasi-separated}；
\item 泛闭；见《空间的态射》引理
\ref{spaces-morphisms-lemma-base-change-universally-closed}
以及《空间上的下降》引理
\ref{spaces-descent-lemma-descending-property-universally-closed}；
\item 泛开；见《空间的态射》引理
\ref{spaces-morphisms-lemma-base-change-universally-open}
以及《空间上的下降》引理
\ref{spaces-descent-lemma-descending-property-universally-open}；
\item 泛商映射；见《空间的态射》引理
\ref{spaces-morphisms-lemma-base-change-universally-submersive}
以及《空间上的下降》引理
\ref{spaces-descent-lemma-descending-property-universally-submersive}；
\item 泛同胚；见《空间的态射》引理
\ref{spaces-morphisms-lemma-base-change-universal-homeomorphism}
以及《空间上的下降》引理
\ref{spaces-descent-lemma-descending-property-universal-homeomorphism}；
\item 满射；见《空间的态射》引理
\ref{spaces-morphisms-lemma-base-change-surjective}
以及《空间上的下降》引理
\ref{spaces-descent-lemma-descending-property-surjective}；
\item 泛单射；见《空间的态射》引理
\ref{spaces-morphisms-lemma-base-change-universally-injective}
以及《空间上的下降》引理
\ref{spaces-descent-lemma-descending-property-universally-injective}；
\item 局部有限型；见《空间的态射》引理
\ref{spaces-morphisms-lemma-base-change-finite-type}
以及《空间上的下降》引理
\ref{spaces-descent-lemma-descending-property-locally-finite-type}；
\item 局部有限表示；见《空间的态射》引理
\ref{spaces-morphisms-lemma-base-change-finite-presentation}
以及《空间上的下降》引理
\ref{spaces-descent-lemma-descending-property-locally-finite-presentation}；
\item 有限型；见《空间的态射》引理
\ref{spaces-morphisms-lemma-base-change-finite-type}
以及《空间上的下降》引理
\ref{spaces-descent-lemma-descending-property-finite-type}；
\item 有限表示；见《空间的态射》引理
\ref{spaces-morphisms-lemma-base-change-finite-presentation}
以及《空间上的下降》引理
\ref{spaces-descent-lemma-descending-property-finite-presentation}；
\item 平坦；见《空间的态射》引理
\ref{spaces-morphisms-lemma-base-change-flat}
以及《空间上的下降》引理
\ref{spaces-descent-lemma-descending-property-flat}；
\item 开浸入；见《空间的态射》第
\ref{spaces-morphisms-section-immersions} 节
以及《空间上的下降》引理
\ref{spaces-descent-lemma-descending-property-open-immersion}；
\item 同构；见《空间上的下降》引理
\ref{spaces-descent-lemma-descending-property-isomorphism}；
\item 仿射；见《空间的态射》引理
\ref{spaces-morphisms-lemma-base-change-affine}
以及《空间上的下降》引理
\ref{spaces-descent-lemma-descending-property-affine}；
\item 闭浸入；见《空间的态射》第
\ref{spaces-morphisms-section-immersions} 节
以及《空间上的下降》引理
\ref{spaces-descent-lemma-descending-property-closed-immersion}；
\item 分离；见《空间的态射》引理
\ref{spaces-morphisms-lemma-base-change-separated}
以及《空间上的下降》引理
\ref{spaces-descent-lemma-descending-property-separated}；
\item 真；见《空间的态射》引理
\ref{spaces-morphisms-lemma-base-change-proper}
以及《空间上的下降》引理
\ref{spaces-descent-lemma-descending-property-proper}；
\item 拟仿射；见《空间的态射》引理
\ref{spaces-morphisms-lemma-base-change-quasi-affine}
以及《空间上的下降》引理
\ref{spaces-descent-lemma-descending-property-quasi-affine}；
\item 整；见《空间的态射》引理
\ref{spaces-morphisms-lemma-base-change-integral}
以及《空间上的下降》引理
\ref{spaces-descent-lemma-descending-property-integral}；
\item 有限；见《空间的态射》引理
\ref{spaces-morphisms-lemma-base-change-integral}
以及《空间上的下降》引理
\ref{spaces-descent-lemma-descending-property-finite}；
\item （局部）拟有限；见《空间的态射》引理
\ref{spaces-morphisms-lemma-base-change-quasi-finite}
以及《空间上的下降》引理
\ref{spaces-descent-lemma-descending-property-quasi-finite}；
\item syntomic；见《空间的态射》引理
\ref{spaces-morphisms-lemma-base-change-syntomic}
以及《空间上的下降》引理
\ref{spaces-descent-lemma-descending-property-syntomic}；
\item 光滑；见《空间的态射》引理
\ref{spaces-morphisms-lemma-base-change-smooth}
以及《空间上的下降》引理
\ref{spaces-descent-lemma-descending-property-smooth}；
\item 非分歧；见《空间的态射》引理
\ref{spaces-morphisms-lemma-base-change-unramified}
以及《空间上的下降》引理
\ref{spaces-descent-lemma-descending-property-unramified}；
\item \'etale；见《空间的态射》引理
\ref{spaces-morphisms-lemma-base-change-etale}
以及《空间上的下降》引理
\ref{spaces-descent-lemma-descending-property-etale}；
\item 有限局部自由；见《空间的态射》引理
\ref{spaces-morphisms-lemma-base-change-finite-locally-free}
以及《空间上的下降》引理
\ref{spaces-descent-lemma-descending-property-finite-locally-free}；
\item 单态射；见《空间的态射》引理
\ref{spaces-morphisms-lemma-base-change-monomorphism}
以及《空间上的下降》引理
\ref{spaces-descent-lemma-descending-property-monomorphism}；
\item 浸入；见《空间的态射》第
\ref{spaces-morphisms-section-immersions} 节
以及《空间上的下降》引理
\ref{spaces-descent-lemma-descending-fppf-property-immersion}；
\item 局部分离；见《空间的态射》引理
\ref{spaces-morphisms-lemma-base-change-separated}
以及《空间上的下降》引理
\ref{spaces-descent-lemma-descending-fppf-property-locally-separated}。
\end{enumerate}

\begin{lemma}
\label{stacks-properties-lemma-property-spaces-too}
设 $P$ 是如上的代数空间态射的性质。
设 $f : \mathcal{X} \to \mathcal{Y}$ 是可由代数空间表示的代数叠态射。
下列条件等价：
\begin{enumerate}
\item $f$ 具有性质 $P$；
\item 对每个代数空间 $Z$ 以及态射 $Z \to \mathcal{Y}$，
态射 $Z \times_\mathcal{Y} \mathcal{X} \to Z$ 具有性质 $P$。
\end{enumerate}
\end{lemma}

\begin{proof}
(2) $\Rightarrow$ (1) 是直接的。设 (1) 成立。
取 (2) 中的 $Z \to \mathcal{Y}$。选取概形 $U$ 以及满的 \'etale 态射
$U \to Z$。由假设，态射
$U \times_\mathcal{Y} \mathcal{X} \to U$ 具有性质 $P$。但图表
$$
\xymatrix{
U \times_\mathcal{Y} \mathcal{X} \ar[d] \ar[r] &
Z \times_\mathcal{Y} \mathcal{X} \ar[d] \\
U \ar[r] & Z
}
$$
是笛卡儿的；由于 $\{U \to Z\}$ 是 fppf 覆盖，右侧竖直箭头因而具有
性质 $P$。
\end{proof}

\noindent
下一个引理说明，只需在沿一个满、平坦且局部有限表示的态射作基变换后
检验性质 $P$。

\begin{lemma}
\label{stacks-properties-lemma-check-property-covering}
设 $P$ 是如上的代数空间态射的性质。
设 $f : \mathcal{X} \to \mathcal{Y}$ 是可由代数空间表示的代数叠态射。
设 $W$ 是代数空间，并设 $W \to \mathcal{Y}$ 满、局部有限表示且平坦。
令 $V = W \times_\mathcal{Y} \mathcal{X}$。则
$$
(f\text{ 具有性质 }P) \Leftrightarrow (\text{投影 }V \to W\text{ 具有性质 }P).
$$
\end{lemma}

\begin{proof}
由引理 \ref{stacks-properties-lemma-property-spaces-too} 可得从左到右的蕴含。
设 $V \to W$ 具有性质 $P$。令 $T$ 为概形，并取态射
$T \to \mathcal{Y}$。考虑代数空间的交换图表
$$
\xymatrix{
T \times_\mathcal{Y} \mathcal{X} \ar[d] &
T \times_\mathcal{Y} V \ar[d] \ar[l] \ar[r] &
V \ar[d] \\
T & T \times_\mathcal{Y} W \ar[l] \ar[r] & W
}
$$
其中两个方块都是笛卡儿的。
左下方态射满、平坦且局部有限表示，故
$\{T \times_\mathcal{Y} V \to T\}$ 是 fppf 覆盖。
因此，右侧竖直箭头具有性质 $P$ 蕴含左侧竖直箭头也具有性质 $P$。
\end{proof}

\begin{lemma}
\label{stacks-properties-lemma-check-property-weak-covering}
设 $P$ 是如上的代数空间态射的性质。
设 $f : \mathcal{X} \to \mathcal{Y}$ 是可由代数空间表示的代数叠态射。
设 $\mathcal{Z} \to \mathcal{Y}$ 是可由代数空间表示、满、平坦且
局部有限表示的代数叠态射。
令 $\mathcal{W} = \mathcal{Z} \times_\mathcal{Y} \mathcal{X}$。则
$$
(f\text{ 具有性质 }P) \Leftrightarrow
(\text{投影 }\mathcal{W} \to \mathcal{Z}\text{ 具有性质 }P).
$$
\end{lemma}

\begin{proof}
选取代数空间 $W$ 以及满、平坦且局部有限表示的态射
$W \to \mathcal{Z}$。由上述讨论，复合
$W \to \mathcal{Y}$ 也满、平坦且局部有限表示。记
$V = W \times_\mathcal{Z} \mathcal{W} = V \times_\mathcal{Y} \mathcal{X}$.
由引理 \ref{stacks-properties-lemma-check-property-covering} 可知，$f$ 具有性质
$\mathcal{P}$ 当且仅当 $V \to W$ 具有该性质，而
$\mathcal{W} \to \mathcal{Z}$ 具有性质 $\mathcal{P}$ 当且仅当
$V \to W$ 具有该性质。引理得证。
\end{proof}

\begin{lemma}
\label{stacks-properties-lemma-check-property-after-precomposing}
\leavevmode\par
设 $P$ 是如上的代数空间态射的性质。
设 $\tau \in \{\etale, smooth, syntomic, fppf\}$。
设 $\mathcal{X} \to \mathcal{Y}$ 和 $\mathcal{Y} \to \mathcal{Z}$
是可由代数空间表示的代数叠态射。假设
\begin{enumerate}
\item $\mathcal{X} \to \mathcal{Y}$ 满，并且是 \'etale、光滑、syntomic，
或者平坦且局部有限表示；
\item 该复合具有性质 $P$；
\item $P$ 在 $\tau$ 拓扑下对源是局部的。
\end{enumerate}
则 $\mathcal{Y} \to \mathcal{Z}$ 具有性质 $P$。
\end{lemma}

\begin{proof}
令 $Z$ 为概形，并取态射 $Z \to \mathcal{Z}$。令
$X = \mathcal{X} \times_\mathcal{Z} Z$、
$Y = \mathcal{Y} \times_\mathcal{Z} Z$。由 (1)，$\{X \to Y\}$
是代数空间的一个 $\tau$ 覆盖；由 (2)，$X \to Z$ 具有性质 $P$。
再由 (3)，这蕴含 $Y \to Z$ 具有性质 $P$，结论得证。
\end{proof}

\begin{lemma}
\label{stacks-properties-lemma-representable-in-terms-presentations}
设 $g : \mathcal{X}' \to \mathcal{X}$ 是可由代数空间表示的代数叠态射。
设 $[U/R] \to \mathcal{X}$ 是一个表示。令
$U' = U \times_\mathcal{X} \mathcal{X}'$、
$R' = R \times_\mathcal{X} \mathcal{X}'$。
则存在形如 $(U', R', s', t', c')$ 的代数空间群胚和一个表示
$[U'/R'] \to \mathcal{X}'$，并且图表
$$
\xymatrix{
[U'/R'] \ar[d]_{[\text{pr}]} \ar[r] & \mathcal{X}' \ar[d]^g \\
[U/R] \ar[r] & \mathcal{X}
}
$$
是 $2$-交换的，其中态射 $[\text{pr}]$ 来自群胚态射
$\text{pr} : (U', R', s', t', c') \to (U, R, s, t, c)$.
\end{lemma}

\begin{proof}
由于 $U \to \mathcal{Y}$ 满且光滑（见《代数叠》引理
\ref{algebraic-lemma-smooth-quotient-smooth-presentation}），
基变换 $U' \to \mathcal{X}'$ 也满且光滑。
因此，由《代数叠》引理 \ref{algebraic-lemma-stack-presentation}，
为得到光滑群胚 $(U', R', s', t', c')$ 和表示
$[U'/R'] \to \mathcal{X}'$，只需证明
$R' = U' \times_{\mathcal{X}'} U'$。
利用 $R = V \times_\mathcal{Y} V$（见《空间中的群胚》引理
\ref{spaces-groupoids-lemma-quotient-stack-2-cartesian}），
这由下式推出：
$$
R' =
U \times_\mathcal{X} U \times_\mathcal{X} \mathcal{X}' =
(U \times_\mathcal{X} \mathcal{X}')
\times_{\mathcal{X}'}
(U \times_\mathcal{X} \mathcal{X}')
$$
见《范畴》引理 \ref{categories-lemma-associativity-2-fibre-product} 和
\ref{categories-lemma-2-fibre-product-erase-factor}。
显然，投影态射 $U' \to U$ 和 $R' \to R$ 给出所需的群胚态射
$\text{pr} : (U', R', s', t', c') \to (U, R, s, t, c)$.
于是商叠的态射 $[\text{pr}]$，由《空间中的群胚》引理
\ref{spaces-groupoids-lemma-quotient-stack-functorial}。

\medskip\noindent
还需证明该图表 $2$-交换。显然，图表
$$
\xymatrix{
U' \ar[d]_{\text{pr}_U} \ar[r]_{f'} & \mathcal{X}' \ar[d]^g \\
U \ar[r]^f & \mathcal{X}
}
$$
$2$-交换，其中 $\text{pr}_U : U' \to U$ 是投影。
由 $R = U \times_\mathcal{X} U$、$t = \text{pr}_0$ 和
$s = \text{pr}_1$，在 $\Mor(R, \mathcal{X})$ 中有典范 $2$-箭头
$\tau : f \circ t \to f \circ s$。利用同构
$R' \to U' \times_{\mathcal{X}'} U'$，类似地得到同构
$\tau' : f' \circ t' \to f' \circ s'$。注意
$g \circ f' \circ t' = f \circ t \circ \text{pr}_R$ 以及
$g \circ f' \circ s' = f \circ s \circ \text{pr}_R$，其中
$\text{pr}_R : R' \to R$ 是投影。因此，可以考察是否有
\begin{equation}
\label{stacks-properties-equation-verify}
\tau \star \text{id}_{\text{pr}_R} = \text{id}_g \star \tau'.
\end{equation}
现作两个断言：(1) 若方程 (\ref{stacks-properties-equation-verify}) 成立，则该图表
$2$-交换；(2) 方程 (\ref{stacks-properties-equation-verify}) 成立。
我们略去两个断言的证明。提示：(1) 来自《代数叠》引理
\ref{algebraic-lemma-map-space-into-stack} 中
$f = f_{can}$ 和 $f' = f'_{can}$ 的构造；(2) 可由仔细展开定义得到。
\end{proof}

\begin{remark}
\label{stacks-properties-remark-representable-over}
设 $\mathcal{Y}$ 是代数叠。考虑如下 $2$-范畴：
\begin{enumerate}
\item 对象是可由代数空间表示的态射
$f : \mathcal{X} \to \mathcal{Y}$；
\item 一个 $1$-态射
$(g, \beta) :
(f_1 : \mathcal{X}_1 \to \mathcal{Y})
\to
(f_2 : \mathcal{X}_2 \to \mathcal{Y})$
由态射 $g : \mathcal{X}_1 \to \mathcal{X}_2$ 和 $2$-态射
$\beta : f_1 \to f_2 \circ g$ 组成；
\item 两个这样的态射
$(g, \beta), (g', \beta') :
(f_1 : \mathcal{X}_1 \to \mathcal{Y})
\to
(f_2 : \mathcal{X}_2 \to \mathcal{Y})$
之间的 $2$-态射，是满足
$(\text{id}_{f_2} \star \alpha) \circ \beta = \beta'$
的 $2$-态射 $\alpha : g \to g'$。
\end{enumerate}
类比《空间上的拓扑》第 \ref{spaces-topologies-section-procedure} 节的记号，
将此 $2$-范畴记作 $\textit{Spaces}/\mathcal{Y}$。
我们断言，在此 $2$-范畴中，态射范畴
$$
\Mor_{\textit{Spaces}/\mathcal{Y}}(
(f_1 : \mathcal{X}_1 \to \mathcal{Y}),
(f_2 : \mathcal{X}_2 \to \mathcal{Y}))
$$
全都是集合胚。事实上，一个 $2$-态射 $\alpha$ 是如下规则：它为
$\mathcal{X}_1$ 的每个对象 $x_1$ 指定相关纤维范畴
$\mathcal{X}_2$ 中的一个同构
$\alpha_{x_1} : g(x_1) \longrightarrow g'(x_1)$
使得图表
$$
\xymatrix{
& f_2(x_1) \ar[ld]_{\beta_{x_1}} \ar[rd]^{\beta'_{x_1}} \\
f_2(g(x_1)) \ar[rr]^{f_2(\alpha_{x_1})} & &
f_2(g'(x_1))
}
$$
交换。但由于 $f_2$ 忠实（见《代数叠》引理
\ref{algebraic-lemma-characterize-representable-by-algebraic-spaces}），
这意味着只要 $\alpha_{x_1}$ 存在，它便是唯一的！换言之，
$2$-范畴 $\textit{Spaces}/\mathcal{Y}$ 与一个范畴非常接近：
若用 $1$-态射的同构类替代 $1$-态射，就得到一个范畴。
我们常常不加说明地作此替换。
\end{remark}




\section{代数叠的点}
\label{stacks-properties-section-points}

\noindent
设 $\mathcal{X}$ 是代数叠，$K, L$ 是两个域，并设
$p : \Spec(K) \to \mathcal{X}$ 和
$q : \Spec(L) \to \mathcal{X}$ 是态射。
若存在域 $\Omega$ 和一个 $2$-交换图表
$$
\xymatrix{
\Spec(\Omega) \ar[r] \ar[d] &
\Spec(L) \ar[d]^q \\
\Spec(K) \ar[r]^p &
\mathcal{X}.
}
$$
则称 $p$ 与 $q$ {\it 等价}。

\begin{lemma}
\label{stacks-properties-lemma-equivalence}
上述概念确实在从域的谱到代数叠 $\mathcal{X}$ 的态射上定义了一个等价关系。
\end{lemma}

\begin{proof}
显然，该关系自反且对称。因此只需证明其传递性。
这归结为如下问题：给定图表
$$
\xymatrix{
\Spec(\Omega) \ar[r]_b \ar[d]_a &
\Spec(L) \ar[d]^q & \Spec(\Omega') \ar[l]^{b'} \ar[d]^{a'} \\
\Spec(K) \ar[r]^p &
\mathcal{X} &
\Spec(K') \ar[l]_{p'}
}
$$
且两个方块都 $2$-交换，需要证明 $p$ 与 $p'$ 等价。
由 $2$-Yoneda 引理（见《代数叠》第
\ref{algebraic-section-2-yoneda} 节），态射 $p$、$p'$ 和 $q$
分别由 $\mathcal{X}$ 在 $\Spec(K)$、$\Spec(K')$ 和 $\Spec(L)$
上纤维范畴中的对象 $x$、$x'$ 和 $y$ 给出。
两个方块的 $2$-交换性意味着存在同构
$\alpha : a^*x \to b^*y$ 和 $\alpha' : (a')^*x' \to (b')^*y$
分别位于 $\mathcal{X}$ 在 $\Spec(\Omega)$ 和 $\Spec(\Omega')$
上的纤维范畴中。任取域 $\Omega''$ 以及在 $L$ 上相同的嵌入
$\Omega \to \Omega''$ 和 $\Omega' \to \Omega''$。
于是可将上图扩充为
$$
\xymatrix{
& \Spec(\Omega'') \ar[ld]_c \ar[d]^{q'} \ar[rd]^{c'} \\
\Spec(\Omega) \ar[r]_b \ar[d]_a &
\Spec(L) \ar[d]^q & \Spec(\Omega') \ar[l]^{b'} \ar[d]^{a'} \\
\Spec(K) \ar[r]^p &
\mathcal{X} &
\Spec(K') \ar[l]_{p'}
}
$$
其中两个三角形交换，并且
$$
(q')^*(\alpha')^{-1} \circ (q')^*\alpha :
(a \circ c)^*x
\longrightarrow
(a' \circ c')^*x'
$$
是 $\Spec(\Omega'')$ 上纤维范畴中的同构。
故 $p$ 与 $p'$ 等价，正如所需。
\end{proof}

\begin{definition}
\label{stacks-properties-definition-points}
设 $\mathcal{X}$ 是代数叠。$\mathcal{X}$ 的一个{\it 点}，是从域的谱到
$\mathcal{X}$ 的态射所组成的一个等价类。
$\mathcal{X}$ 的点集记作 $|\mathcal{X}|$。
\end{definition}

\noindent
这与代数空间之点的定义一致；见《空间的性质》定义
\ref{spaces-properties-definition-points}。
此外，对概形而言，这恢复了通常的点的概念；见《空间的性质》引理
\ref{spaces-properties-lemma-scheme-points}。
若 $f : \mathcal{X} \to \mathcal{Y}$ 是代数叠态射，则它诱导映射
$|f| : |\mathcal{X}| \to |\mathcal{Y}|$，将代表元
$x : \Spec(K) \to \mathcal{X}$ 映为代表元
$f \circ x : \Spec(K) \to \mathcal{Y}$。该映射适定：
$2$-同构的 $1$-态射与一个 $1$-态射前复合或后复合后仍 $2$-同构，
因为可以与所给 $1$-态射的恒等态射作水平前复合或后复合。
这在任意（严格）$(2, 1)$-范畴中都成立。若
$$
\xymatrix{
\mathcal{X} \ar[d] \ar[r] & \mathcal{Y} \ar[d] \\
\mathcal{W} \ar[r] & \mathcal{Z}
}
$$
是代数叠的 $2$-交换图表，则集合图表
$$
\xymatrix{
|\mathcal{X}| \ar[d] \ar[r] & |\mathcal{Y}| \ar[d] \\
|\mathcal{W}| \ar[r] & |\mathcal{Z}|
}
$$
交换。特别地，若 $\mathcal{X} \to \mathcal{Y}$ 是等价，
则 $|\mathcal{X}| \to |\mathcal{Y}|$ 是双射。

\begin{lemma}
\label{stacks-properties-lemma-points-cartesian}
设
$$
\xymatrix{
\mathcal{Z} \times_\mathcal{Y} \mathcal{X} \ar[r] \ar[d] &
\mathcal{X} \ar[d] \\
\mathcal{Z} \ar[r] & \mathcal{Y}
}
$$
是代数叠的纤维积。则点集映射
$$
|\mathcal{Z} \times_\mathcal{Y} \mathcal{X}|
\longrightarrow
|\mathcal{Z}| \times_{|\mathcal{Y}|} |\mathcal{X}|
$$
是满射。
\end{lemma}

\begin{proof}
事实上，设给定域 $K$、$L$ 和态射
$\Spec(K) \to \mathcal{X}$、$\Spec(L) \to \mathcal{Z}$。
它们作为 $|\mathcal{Y}|$ 的元素相同这一假设，意味着存在公共扩张
$M/K$ 和 $M/L$，使得
$\Spec(M) \to \Spec(K) \to \mathcal{X} \to \mathcal{Y}$ 和
$\Spec(M) \to \Spec(L) \to \mathcal{Z} \to \mathcal{Y}$
是 $2$-同构的。这恰好就是得到态射
$\Spec(M) \to \mathcal{Z} \times_\mathcal{Y} \mathcal{X}$ 的条件。
\end{proof}

\begin{lemma}
\label{stacks-properties-lemma-characterize-surjective}
设 $f : \mathcal{X} \to \mathcal{Y}$ 是可由代数空间表示的代数叠态射。
下列条件等价：
\begin{enumerate}
\item $|f| : |\mathcal{X}| \to |\mathcal{Y}|$ 是满射；
\item $f$ 是满射（取第 \ref{stacks-properties-section-properties-morphisms} 节的含义）。
\end{enumerate}
\end{lemma}

\begin{proof}
设 (1) 成立。令 $T \to \mathcal{Y}$ 为源是概形的态射。
为证明 (2)，需证明代数空间态射
$T \times_\mathcal{Y} \mathcal{X} \to T$ 是满射。
由《空间的态射》定义 \ref{spaces-morphisms-definition-surjective}，
这意味着要证明 $|T \times_\mathcal{Y} \mathcal{X}| \to |T|$ 是满射。
应用引理 \ref{stacks-properties-lemma-points-cartesian}，可见这由 (1) 推出。

\medskip\noindent
反之，设 (2) 成立。令 $y : \Spec(K) \to \mathcal{Y}$
为从域的谱到 $\mathcal{Y}$ 的态射。由假设，代数空间态射
$\Spec(K) \times_{y, \mathcal{Y}} \mathcal{X} \to \Spec(K)$
是满射。由《空间的态射》定义
\ref{spaces-morphisms-definition-surjective}，这意味着存在域扩张
$K'/K$ 和态射
$\Spec(K') \to \Spec(K) \times_{y, \mathcal{Y}} \mathcal{X}$
使下图中的左方块交换：
$$
\xymatrix{
\Spec(K') \ar[r] \ar[d] &
\Spec(K) \times_{y, \mathcal{Y}} \mathcal{X} \ar[d] \ar[r] &
\mathcal{X} \ar[d]
\\
\Spec(K) \ar@{=}[r] &
\Spec(K) \ar[r]^-y &
\mathcal{Y}
}
$$
这说明 $|X| \to |\mathcal{Y}|$ 是满射。
\end{proof}

\noindent
下面的引理说明如何利用一个表示计算点集。

\begin{lemma}
\label{stacks-properties-lemma-points-presentation}
设 $\mathcal{X}$ 是代数叠。设 $\mathcal{X} = [U/R]$
是 $\mathcal{X}$ 的一个表示；见《代数叠》定义
\ref{algebraic-definition-presentation}。
则 $|R| \to |U| \times |U|$ 的像是一个等价关系，
且 $|\mathcal{X}|$ 是 $|U|$ 关于该等价关系的商。
\end{lemma}

\begin{proof}
该假设意味着我们有代数空间中的一个光滑群胚
$(U, R, s, t, c)$，以及等价
$f : [U/R] \to \mathcal{X}$。可以假设 $\mathcal{X} = [U/R]$。
诱导态射 $p : U \to \mathcal{X}$ 光滑且满；见《代数叠》引理
\ref{algebraic-lemma-smooth-quotient-smooth-presentation}。
因此，由引理 \ref{stacks-properties-lemma-characterize-surjective}，
$|U| \to |\mathcal{X}|$ 是满射。注意
$R = U \times_\mathcal{X} U$；见《空间中的群胚》引理
\ref{spaces-groupoids-lemma-quotient-stack-2-cartesian}。
故引理 \ref{stacks-properties-lemma-points-cartesian} 蕴含映射
$$
|R| \longrightarrow |U| \times_{|\mathcal{X}|} |U|
$$
是满射。因此，$|R| \to |U| \times |U|$ 的像恰好由所有满足
$u_1$ 与 $u_2$ 在 $|\mathcal{X}|$ 中像相同的偶
$(u_1, u_2) \in |U| \times |U|$ 构成。
合并这两个陈述即得引理。
\end{proof}

\begin{remark}
\label{stacks-properties-remark-more-general-presentation}
引理 \ref{stacks-properties-lemma-points-presentation} 的结果可作如下推广。
设 $\mathcal{X}$ 是代数叠。设 $U$ 是代数空间，并设
$f : U \to \mathcal{X}$ 是满射（由第
\ref{stacks-properties-section-properties-morphisms} 节，这一说法有意义）。
令 $R = U \times_\mathcal{X} U$，令 $(U, R, s, t, c)$
为代数空间中的群胚，并令 $f_{can} : [U/R] \to \mathcal{X}$
为《代数叠》引理 \ref{algebraic-lemma-map-space-into-stack}
中构造的典范态射。
则 $|R| \to |U| \times |U|$ 的像是一个等价关系，且
$|\mathcal{X}| = |U|/|R|$。引理 \ref{stacks-properties-lemma-points-presentation}
的证明无需改动便可适用。（当然，一般而言 $[U/R]$ 不是代数叠，
且一般而言 $f_{can}$ 不是同构。）
\end{remark}

\begin{lemma}
\label{stacks-properties-lemma-topology-points}
在代数叠的点集上，存在唯一的拓扑满足下列性质：
\begin{enumerate}
\item 对每个代数叠态射 $\mathcal{X} \to \mathcal{Y}$，
映射 $|\mathcal{X}| \to |\mathcal{Y}|$ 连续；
\item 对每个平坦且局部有限表示的态射 $U \to \mathcal{X}$，
其中 $U$ 是代数空间，拓扑空间映射
$|U| \to |\mathcal{X}|$ 连续且开。
\end{enumerate}
\end{lemma}

\begin{proof}
选取一个满、平坦且局部有限表示的态射
$p : U \to \mathcal{X}$，其中 $U$ 是代数空间。
由代数叠的定义，这样的态射存在，因为光滑态射平坦且局部有限表示
（见《空间的态射》引理
\ref{spaces-morphisms-lemma-smooth-locally-finite-presentation} 和
\ref{spaces-morphisms-lemma-smooth-flat}）。
按如下规则在 $|\mathcal{X}|$ 上定义拓扑：
$W \subset |\mathcal{X}|$ 是开集，当且仅当
$|p|^{-1}(W)$ 在 $|U|$ 中是开集。
为证明这与 $p$ 的选择无关，令 $p' : U' \to \mathcal{X}$
为另一个从代数空间到 $\mathcal{X}$ 的满、平坦且局部有限表示的态射。
令 $U'' = U \times_\mathcal{X} U'$，于是有 $2$-交换图表
$$
\xymatrix{
U'' \ar[r] \ar[d] & U' \ar[d] \\
U \ar[r] & \mathcal{X}
}
$$
由于 $U \to \mathcal{X}$ 和 $U' \to \mathcal{X}$ 满、平坦且
局部有限表示，由引理 \ref{stacks-properties-lemma-property-spaces-too} 可知
$U'' \to U'$ 和 $U'' \to U$ 也满、平坦且局部有限表示。
因此，映射 $|U''| \to |U'|$ 和 $|U''| \to |U|$
连续、开且满；见《空间的态射》定义
\ref{spaces-morphisms-definition-surjective} 和引理
\ref{spaces-morphisms-lemma-fppf-open}。
这显然说明我们的定义与 $p : U \to \mathcal{X}$ 的选择无关。

\medskip\noindent
设 $f : \mathcal{X} \to \mathcal{Y}$ 是代数叠态射。
由《代数叠》引理 \ref{algebraic-lemma-lift-morphism-presentations}，
可以找到一个 $2$-交换图表
$$
\xymatrix{
U \ar[d]_x \ar[r]_a & V \ar[d]^y \\
\mathcal{X} \ar[r]^f & \mathcal{Y}
}
$$
其中竖直箭头满且光滑。考虑相应的集合交换图表
$$
\xymatrix{
|U| \ar[d]_{|x|} \ar[r]_{|a|} & |V| \ar[d]^{|y|} \\
|\mathcal{X}| \ar[r]^{|f|} & |\mathcal{Y}|
}
$$
若 $W \subset |\mathcal{Y}|$ 是开集，则由上述定义，这恰好意味着
$|y|^{-1}(W)$ 在 $|V|$ 中是开集。由于 $|a|$ 连续，可得
$|a|^{-1}|y|^{-1}(W) = |x|^{-1}|f|^{-1}(W)$ 在 $|W|$ 中是开集；
按定义，这意味着 $|f|^{-1}(W)$ 在 $|\mathcal{X}|$ 中是开集。
故 $|f|$ 连续。

\medskip\noindent
最后，需要证明：若 $U$ 是代数空间，且
$U \to \mathcal{X}$ 平坦并局部有限表示，则
$|U| \to |\mathcal{X}|$ 是开映射。令 $V \to \mathcal{X}$
为满、平坦且局部有限表示的态射，其中 $V$ 是代数空间。
考虑交换图表
$$
\xymatrix{
|U \times_\mathcal{X} V| \ar[r]_e \ar[rd]_f &
|U| \times_{|\mathcal{X}|} |V| \ar[d]_c \ar[r]_d &
|V| \ar[d]^b \\
& |U| \ar[r]^a & |\mathcal{X}|
}
$$
态射 $U \times_\mathcal{X} V \to U$ 是满射，亦即
$f : |U \times_\mathcal{X} V| \to |U|$ 是满射。
左上方的水平箭头是满射；见引理 \ref{stacks-properties-lemma-points-cartesian}。
态射 $U \times_\mathcal{X} V \to V$ 平坦且局部有限表示，
故 $d \circ e : |U \times_\mathcal{X} V| \to |V|$ 是开映射；
见《空间的态射》引理 \ref{spaces-morphisms-lemma-fppf-open}。
任取开集 $W \subset |U|$。上述性质蕴含
$b^{-1}(a(W)) = (d \circ e)(f^{-1}(W))$ 是开集；
按构造，这意味着 $a(W)$ 是开集，正如所需。
\end{proof}

\begin{definition}
\label{stacks-properties-definition-topological-space}
设 $\mathcal{X}$ 是代数叠。$\mathcal{X}$ 的底层{\it 拓扑空间}，
是赋予引理 \ref{stacks-properties-lemma-topology-points} 中所构造拓扑的点集
$|\mathcal{X}|$。
\end{definition}

\noindent
若 $\mathcal{X}$ 是代数空间，则此定义与
$|\mathcal{X}|$ 上已有的拓扑一致。

\begin{lemma}
\label{stacks-properties-lemma-space-locally-quasi-compact}
设 $\mathcal{X}$ 是代数叠。$|\mathcal{X}|$ 的每个点都有一个
由拟紧开邻域组成的基本邻域系。
特别地，$|\mathcal{X}|$ 在《拓扑》定义
\ref{topology-definition-locally-quasi-compact} 的意义下局部拟紧。
\end{lemma}

\begin{proof}
这形式地来自如下事实：存在概形 $U$ 以及拓扑空间的满、开、连续映射
$U \to |\mathcal{X}|$。事实上，若 $U \to \mathcal{X}$ 满且光滑，
则引理 \ref{stacks-properties-lemma-topology-points} 保证
$|U| \to |\mathcal{X}|$ 连续、满且开。
\end{proof}











\section{满态射}
\label{stacks-properties-section-surjective}

\noindent
设 $f : \mathcal{X} \to \mathcal{Y}$ 是可由代数空间表示的代数叠态射。
我们已经在第 \ref{stacks-properties-section-properties-morphisms} 节定义了
$f$ 为满射的含义。引理 \ref{stacks-properties-lemma-characterize-surjective}
说明，这等价于要求
$|f| : |\mathcal{X}| \to |\mathcal{Y}|$ 是满射。
因此可以给出如下定义。

\begin{definition}
\label{stacks-properties-definition-surjective}
设 $f : \mathcal{X} \to \mathcal{Y}$ 是代数叠态射。
若相应拓扑空间的映射
$|f| : |\mathcal{X}| \to |\mathcal{Y}|$ 是满射，
则称 $f$ {\it 满}。
\end{definition}

\noindent
下面给出一些引理。

\begin{lemma}
\label{stacks-properties-lemma-composition-surjective}
满态射的复合是满态射。
\end{lemma}

\begin{proof}
略。
\end{proof}

\begin{lemma}
\label{stacks-properties-lemma-base-change-surjective}
满态射的基变换是满态射。
\end{lemma}

\begin{proof}
略。提示：使用引理 \ref{stacks-properties-lemma-points-cartesian}。
\end{proof}

\begin{lemma}
\label{stacks-properties-lemma-descent-surjective}
设 $f : \mathcal{X} \to \mathcal{Y}$ 是代数叠态射。
设 $\mathcal{Y}' \to \mathcal{Y}$ 是代数叠的满态射。
若 $f$ 的基变换
$f' : \mathcal{Y}' \times_\mathcal{Y} \mathcal{X}
\to \mathcal{Y}'$ 是满射，则 $f$ 是满射。
\end{lemma}

\begin{proof}
由引理 \ref{stacks-properties-lemma-points-cartesian} 立即可得。
\end{proof}

\begin{lemma}
\label{stacks-properties-lemma-surjective-permanence}
设 $\mathcal{X} \to \mathcal{Y} \to \mathcal{Z}$ 是代数叠态射。
若 $\mathcal{X} \to \mathcal{Z}$ 是满射，则
$\mathcal{Y} \to \mathcal{Z}$ 也是满射。
\end{lemma}

\begin{proof}
立即可得。
\end{proof}












\section{拟紧代数叠}
\label{stacks-properties-section-quasi-compact}

\noindent
由《空间的性质》引理
\ref{spaces-properties-lemma-quasi-compact-space}，
下述定义与代数空间的相应定义等价。

\begin{definition}
\label{stacks-properties-definition-quasi-compact}
设 $\mathcal{X}$ 是代数叠。当且仅当
$|\mathcal{X}|$ 拟紧时，称 $\mathcal{X}$ {\it 拟紧}。
\end{definition}

\begin{lemma}
\label{stacks-properties-lemma-quasi-compact-stack}
设 $\mathcal{X}$ 是代数叠。下列条件等价：
\begin{enumerate}
\item $\mathcal{X}$ 拟紧；
\item 存在满的光滑态射 $U \to \mathcal{X}$，其中 $U$ 是仿射概形；
\item 存在满的光滑态射 $U \to \mathcal{X}$，其中 $U$ 是拟紧概形；
\item 存在满的光滑态射 $U \to \mathcal{X}$，其中 $U$ 是拟紧代数空间；
\item 存在代数叠的满态射 $\mathcal{U} \to \mathcal{X}$，
其中 $\mathcal{U}$ 拟紧。
\end{enumerate}
\end{lemma}

\begin{proof}
我们将使用引理 \ref{stacks-properties-lemma-characterize-surjective}。
设 $\mathcal{U}$ 和 $\mathcal{U} \to \mathcal{X}$ 如 (5)。
由于 $|\mathcal{U}| \to |\mathcal{X}|$ 满且连续，
可知 $|\mathcal{X}|$ 拟紧。因此 (5) 蕴含 (1)。蕴含
(2) $\Rightarrow$ (3) $\Rightarrow$ (4) $\Rightarrow$ (5)
是直接的。设 (1) 成立，即 $\mathcal{X}$ 拟紧，也即
$|\mathcal{X}|$ 拟紧。选取概形 $U$ 以及满的光滑态射
$U \to \mathcal{X}$。由于 $|U| \to |\mathcal{X}|$ 是开映射，
存在拟紧开子集 $U' \subset U$，使得 $|U'| \to |X|$ 满
（并且仍然光滑）。选取有限仿射开覆盖
$U' = U_1 \cup \ldots \cup U_n$。则
$U_1 \amalg \ldots \amalg U_n \to \mathcal{X}$
是满的光滑态射，且其源是仿射概形（见《概形》引理
\ref{schemes-lemma-disjoint-union-affines}）。故 (2) 成立。
\end{proof}

\begin{lemma}
\label{stacks-properties-lemma-finite-disjoint-quasi-compact}
拟紧代数叠的有限不交并是拟紧代数叠。
\end{lemma}

\begin{proof}
这由相应的拓扑事实立即可得。
\end{proof}


\section{由概形性质定义的代数叠性质}
\label{stacks-properties-section-types-properties}

\noindent
概形的任一光滑局部性质都通过下述引理给出代数叠的相应性质。
注意，概形的光滑局部性质也是 \'etale 局部性质，
因为任一 \'etale 覆盖也是光滑覆盖。
因此，对概形的光滑局部性质 $P$，我们知道代数空间具有 $P$
的含义；见《空间的性质》第
\ref{spaces-properties-section-types-properties} 节。

\begin{lemma}
\label{stacks-properties-lemma-type-property}
设 $\mathcal{P}$ 是在光滑拓扑下局部的概形性质；见《下降》定义
\ref{descent-definition-property-local}。
设 $\mathcal{X}$ 是代数叠。下列条件等价：
\begin{enumerate}
\item 对某个概形 $U$ 和某个满的光滑态射 $U \to \mathcal{X}$，
概形 $U$ 具有性质 $\mathcal{P}$；
\item 对每个概形 $U$ 和每个光滑态射 $U \to \mathcal{X}$，
概形 $U$ 具有性质 $\mathcal{P}$；
\item 对某个代数空间 $U$ 和某个满的光滑态射 $U \to \mathcal{X}$，
代数空间 $U$ 具有性质 $\mathcal{P}$；
\item 对每个代数空间 $U$ 和每个光滑态射 $U \to \mathcal{X}$，
代数空间 $U$ 具有性质 $\mathcal{P}$。
\end{enumerate}
若 $\mathcal{X}$ 是概形 $U$，则这等价于 $\mathcal{P}(U)$。
若 $\mathcal{X}$ 是代数空间 $X$，则这等价于
$X$ 具有性质 $\mathcal{P}$。
\end{lemma}

\begin{proof}
设 $U \to \mathcal{X}$ 满且光滑，其中 $U$ 是代数空间。
设 $V \to \mathcal{X}$ 是光滑态射，其中 $V$ 是代数空间。
选取概形 $U'$、$V'$ 以及满的 \'etale 态射
$U' \to U$、$V' \to V$。最后，选取概形 $W$ 以及满的 \'etale 态射
$W \to V' \times_\mathcal{X} U'$。
由于 $W \to V'$ 和 $W \to U'$ 是代数空间的 \'etale 态射与
光滑态射的复合，它们是概形的光滑态射；见《空间的态射》引理
\ref{spaces-morphisms-lemma-etale-smooth} 和
\ref{spaces-morphisms-lemma-composition-smooth}。
此外，由于 $U' \to \mathcal{X}$ 满，$W \to V'$ 也满。因此
$$
\mathcal{P}(U) \Leftrightarrow
\mathcal{P}(U') \Rightarrow
\mathcal{P}(W) \Rightarrow
\mathcal{P}(V') \Leftrightarrow \mathcal{P}(V)
$$
其中两个等价来自代数空间的性质 $\mathcal{P}$ 之定义，
而两个蕴含来自《下降》定义
\ref{descent-definition-property-local}。
这证明了 (3) $\Rightarrow$ (4)。

\medskip\noindent
蕴含 (2) $\Rightarrow$ (1)、(1) $\Rightarrow$ (3)
以及 (4) $\Rightarrow$ (2) 都是直接的。
\end{proof}

\begin{definition}
\label{stacks-properties-definition-type-property}
设 $\mathcal{X}$ 是代数叠，$\mathcal{P}$ 是在光滑拓扑下局部的
概形性质。若引理 \ref{stacks-properties-lemma-type-property} 中任一等价条件成立，
则称 $\mathcal{X}$ {\it 具有性质 $\mathcal{P}$}。
\end{definition}

\begin{remark}
\label{stacks-properties-remark-list-properties-local-smooth-topology}
下面列出一些在光滑拓扑下局部的性质
（注意 fpqc、fppf 和 syntomic 拓扑都强于光滑拓扑）：
\begin{enumerate}
\item 局部 Noether；见《下降》引理
\ref{descent-lemma-Noetherian-local-fppf}；
\item Jacobson；见《下降》引理
\ref{descent-lemma-Jacobson-local-fppf}；
\item 局部 Noether 且满足 $(S_k)$；见《下降》引理
\ref{descent-lemma-Sk-local-syntomic}；
\item Cohen--Macaulay；见《下降》引理
\ref{descent-lemma-CM-local-syntomic}；
\item 既约；见《下降》引理
\ref{descent-lemma-reduced-local-smooth}；
\item 正规；见《下降》引理
\ref{descent-lemma-normal-local-smooth}；
\item 局部 Noether 且满足 $(R_k)$；见《下降》引理
\ref{descent-lemma-Rk-local-smooth}；
\item 正则；见《下降》引理
\ref{descent-lemma-regular-local-smooth}；
\item Nagata；见《下降》引理
\ref{descent-lemma-Nagata-local-smooth}。
\end{enumerate}
\end{remark}

\noindent
概形芽的任一光滑局部性质都给出代数叠的相应性质。
注意，芽的光滑局部性质也是 \'etale 局部性质。
因此，对概形芽的光滑局部性质 $P$，我们知道代数空间 $X$
在 $x \in |X|$ 处具有性质 $P$ 的含义；见《空间的性质》第
\ref{stacks-properties-section-types-properties} 节。

\begin{lemma}
\label{stacks-properties-lemma-local-source-target-at-point}
设 $\mathcal{X}$ 是代数叠，$x \in |\mathcal{X}|$ 是 $\mathcal{X}$ 的点。
设 $\mathcal{P}$ 是概形芽的光滑局部性质；见《下降》定义
\ref{descent-definition-local-at-point}。下列条件等价：
\begin{enumerate}
\item 对任一光滑态射 $U \to \mathcal{X}$，其中 $U$ 是概形，
以及任一满足 $a(u) = x$ 的 $u \in U$，都有 $\mathcal{P}(U, u)$；
\item 对某个光滑态射 $U \to \mathcal{X}$，其中 $U$ 是概形，
以及某个满足 $a(u) = x$ 的 $u \in U$，有 $\mathcal{P}(U, u)$；
\item 对任一光滑态射 $U \to \mathcal{X}$，其中 $U$ 是代数空间，
以及任一满足 $a(u) = x$ 的 $u \in |U|$，代数空间 $U$
在 $u$ 处具有性质 $\mathcal{P}$；
\item 对某个光滑态射 $U \to \mathcal{X}$，其中 $U$ 是代数空间，
以及某个满足 $a(u) = x$ 的 $u \in |U|$，代数空间 $U$
在 $u$ 处具有性质 $\mathcal{P}$。
\end{enumerate}
若 $\mathcal{X}$ 可表示，则这等价于
$\mathcal{P}(\mathcal{X}, x)$。若 $\mathcal{X}$ 是代数空间，
则这等价于 $\mathcal{X}$ 在 $x$ 处具有性质 $\mathcal{P}$。
\end{lemma}

\begin{proof}
取 (3) 中的 $a : U \to \mathcal{X}$ 和 $u \in |U|$。
再取另一个光滑态射 $b : V \to \mathcal{X}$，其中 $V$ 是代数空间，
并取满足 $b(v) = x$ 的 $v \in |V|$。
选取概形 $U'$、\'etale 态射 $U' \to U$ 以及映到 $u$ 的
$u' \in U'$。选取概形 $V'$、\'etale 态射 $V' \to V$
以及映到 $v$ 的 $v' \in V'$。
由引理 \ref{stacks-properties-lemma-points-cartesian}，存在点
$\overline{w} \in |V' \times_\mathcal{X} U'|$ 映到 $u'$ 和 $v'$。
选取概形 $W$ 以及满的 \'etale 态射
$W \to V' \times_\mathcal{X} U'$。可以选取映到 $\overline{w}$ 的
$w \in |W|$（见《空间的性质》引理
\ref{spaces-properties-lemma-characterize-surjective}）。
由于 $W \to V'$ 和 $W \to U'$ 是代数空间的 \'etale 态射与
光滑态射的复合，它们是概形的光滑态射；见《空间的态射》引理
\ref{spaces-morphisms-lemma-etale-smooth} 和
\ref{spaces-morphisms-lemma-composition-smooth}。因此
$$
\mathcal{P}(U, u)
\Leftrightarrow
\mathcal{P}(U', u')
\Leftrightarrow
\mathcal{P}(W, w)
\Leftrightarrow
\mathcal{P}(V', v')
\Leftrightarrow
\mathcal{P}(V, v)
$$
外侧两个等价来自《空间的性质》定义
\ref{spaces-properties-definition-property-at-point}，
另外两个等价来自概形芽的光滑局部性质之含义。
这证明了 (4) $\Rightarrow$ (3)。

\medskip\noindent
蕴含 (1) $\Rightarrow$ (2)、(2) $\Rightarrow$ (4)
以及 (3) $\Rightarrow$ (1) 都是直接的。
\end{proof}

\begin{definition}
\label{stacks-properties-definition-property-at-point}
设 $\mathcal{P}$ 是概形芽的光滑局部性质，$\mathcal{X}$ 是代数叠，
且 $x \in |\mathcal{X}|$。若引理
\ref{stacks-properties-lemma-local-source-target-at-point} 的任一等价条件成立，
则称 $\mathcal{X}$ {\it 在 $x$ 处具有性质 $\mathcal{P}$}。
\end{definition}










\section{代数叠的单态射}
\label{stacks-properties-section-monomorphisms}

\noindent
我们如下定义代数叠的单态射。引理 \ref{stacks-properties-lemma-monomorphism}
将说明，该定义与相应的 $2$-范畴论概念相容。

\begin{definition}
\label{stacks-properties-definition-monomorphism}
设 $f : \mathcal{X} \to \mathcal{Y}$ 是代数叠态射。
若它可由代数空间表示，并且是第
\ref{stacks-properties-section-properties-morphisms} 节意义下的单态射，
则称 $f$ 是{\it 单态射}。
\end{definition}

\noindent
先给出一些基本引理。

\begin{lemma}
\label{stacks-properties-lemma-base-change-monomorphism}
设 $\mathcal{X} \to \mathcal{Y}$ 是代数叠态射，
且 $\mathcal{Z} \to \mathcal{Y}$ 是单态射。
则 $\mathcal{Z} \times_\mathcal{Y} \mathcal{X} \to \mathcal{X}$
是单态射。
\end{lemma}

\begin{proof}
这由第 \ref{stacks-properties-section-properties-morphisms} 节的一般讨论推出。
\end{proof}

\begin{lemma}
\label{stacks-properties-lemma-composition-monomorphism}
代数叠单态射的复合是单态射。
\end{lemma}

\begin{proof}
这由第 \ref{stacks-properties-section-properties-morphisms} 节的一般讨论以及
《空间的态射》引理
\ref{spaces-morphisms-lemma-composition-monomorphism} 推出。
\end{proof}

\begin{lemma}
\label{stacks-properties-lemma-monomorphism}
设 $f : \mathcal{X} \to \mathcal{Y}$ 是代数叠态射。下列条件等价：
\begin{enumerate}
\item $f$ 是单态射；
\item $f$ 全忠实；
\item 对角态射
$\Delta_f : \mathcal{X} \to \mathcal{X} \times_\mathcal{Y} \mathcal{X}$
是等价；
\item 存在代数空间 $W$ 以及满、平坦且局部有限表示的态射
$W \to \mathcal{Y}$，使得
$V = \mathcal{X} \times_\mathcal{Y} W$ 是代数空间，并且
$V \to W$ 是代数空间的单态射。
\end{enumerate}
\end{lemma}

\begin{proof}
(1) 与 (4) 的等价性来自第 \ref{stacks-properties-section-properties-morphisms} 节的
一般讨论，特别是引理 \ref{stacks-properties-lemma-check-representable-covering}
和 \ref{stacks-properties-lemma-check-property-covering}。

\medskip\noindent
(2) 与 (3) 的等价性是《范畴》引理
\ref{categories-lemma-fully-faithful-diagonal-equivalence}。

\medskip\noindent
设等价条件 (2) 和 (3) 成立。由《代数叠》引理
\ref{algebraic-lemma-characterize-representable-by-algebraic-spaces}，
$f$ 可由代数空间表示。此外，$2$-Yoneda 引理与全忠实性共同蕴含：
对每个概形 $T$，函子
$$
\Mor(T, \mathcal{X})
\longrightarrow
\Mor(T, \mathcal{Y})
$$
全忠实。因此，给定态射 $y : T \to \mathcal{Y}$，
在唯一 $2$-同构的意义下，至多存在一个满足
$y \cong f \circ x$ 的态射 $x : T \to \mathcal{X}$。
特别地，给定概形态射 $h : T' \to T$，至多存在一个 $h$ 的提升
$\tilde h : T' \to T \times_\mathcal{Y} \mathcal{X}$。
故 $T \times_\mathcal{Y} \mathcal{X} \to T$ 是代数空间的单态射，
这证明了 (1)。

\medskip\noindent
最后，设 (1) 成立。则对任一概形 $T$ 和态射
$y : T \to \mathcal{Y}$，纤维积
$T \times_\mathcal{Y} \mathcal{X}$ 是代数空间，并且
$T \times_\mathcal{Y} \mathcal{X} \to T$ 是单态射。
因此，在唯一同构的意义下，恰好存在一个偶 $(x, \alpha)$，
其中 $x : T \to \mathcal{X}$ 是态射，而
$\alpha : f \circ x \to y$ 是 $2$-态射。
应用 $2$-Yoneda 引理，这恰好说明 $f$ 全忠实，即 (2) 成立。
\end{proof}

\begin{lemma}
\label{stacks-properties-lemma-monomorphism-injective-points}
\begin{slogan}
叠的单态射在点上是单射。
\end{slogan}
代数叠的单态射诱导点集之间的单射。
\end{lemma}

\begin{proof}
设 $f : \mathcal{X} \to \mathcal{Y}$ 是代数叠的单态射。
设 $x_i : \Spec(K_i) \to \mathcal{X}$ 是态射，并且
$f \circ x_1$ 与 $f \circ x_2$ 定义 $|\mathcal{Y}|$ 的同一元素。
应用定义，可找到公共扩张 $\Omega$、相应态射
$c_i : \Spec(\Omega) \to \Spec(K_i)$，以及 $2$-同构
$\beta : f \circ x_1 \circ c_1 \to f \circ x_1 \circ c_2$。
由于 $f$ 全忠实（见引理 \ref{stacks-properties-lemma-monomorphism}），
可以将 $\beta$ 提升为同构
$\alpha : x_1 \circ c_1 \to x_1 \circ c_2$。
故 $x_1$ 与 $x_2$ 定义 $|\mathcal{X}|$ 的同一点，正如所需。
\end{proof}

\begin{lemma}
\label{stacks-properties-lemma-monomorphism-diagonal}
设 $\mathcal{X} \to \mathcal{X}' \to \mathcal{Y}$ 是代数叠态射。
若 $\mathcal{X} \to \mathcal{X}'$ 是单态射，则典范图表
$$
\xymatrix{
\mathcal{X} \ar[r] \ar[d] &
\mathcal{X} \times_\mathcal{Y} \mathcal{X} \ar[d] \\
\mathcal{X}' \ar[r] &
\mathcal{X}' \times_\mathcal{Y} \mathcal{X}'
}
$$
是纤维积方块。
\end{lemma}

\begin{proof}
由引理 \ref{stacks-properties-lemma-monomorphism}，有
$\mathcal{X} = \mathcal{X} \times_{\mathcal{X}'} \mathcal{X}$。
因此，应用《范畴》引理
\ref{categories-lemma-fibre-product-after-map} 即得结论。
\end{proof}







\section{代数叠的浸入}
\label{stacks-properties-section-immersions}

\noindent
代数叠的浸入定义如下。

\begin{definition}
\label{stacks-properties-definition-immersion}
浸入。
\begin{enumerate}
\item 若一个代数叠态射可表示，并且是第
\ref{stacks-properties-section-properties-morphisms} 节意义下的开浸入，
则称它为{\it 开浸入}。
\item 若一个代数叠态射可表示，并且是第
\ref{stacks-properties-section-properties-morphisms} 节意义下的闭浸入，
则称它为{\it 闭浸入}。
\item 若一个代数叠态射可表示，并且是第
\ref{stacks-properties-section-properties-morphisms} 节意义下的浸入，
则称它为{\it 浸入}。
\end{enumerate}
\end{definition}

\noindent
对我们而言，这并不是理解浸入的最方便方式。更方便的看法是：
浸入是可由代数空间表示、并且是第
\ref{stacks-properties-section-properties-morphisms} 节意义下浸入的代数叠态射。
闭浸入和开浸入亦同。由于这显然等价于刚定义的概念，
以后将不加说明地使用这一刻画。
下面证明关于该概念的几个简单引理。

\begin{lemma}
\label{stacks-properties-lemma-base-change-immersion}
设 $\mathcal{X} \to \mathcal{Y}$ 是代数叠态射，
且 $\mathcal{Z} \to \mathcal{Y}$ 是（闭浸入，分别地，开浸入）。
则 $\mathcal{Z} \times_\mathcal{Y} \mathcal{X} \to \mathcal{X}$
是（闭浸入，分别地，开浸入）。
\end{lemma}

\begin{proof}
这由第 \ref{stacks-properties-section-properties-morphisms} 节的一般讨论推出。
\end{proof}

\begin{lemma}
\label{stacks-properties-lemma-composition-immersion}
代数叠浸入的复合是浸入。闭浸入和开浸入也分别如此。
\end{lemma}

\begin{proof}
这由第 \ref{stacks-properties-section-properties-morphisms} 节的一般讨论以及
《空间》引理 \ref{spaces-lemma-composition-immersions} 推出。
\end{proof}

\begin{lemma}
\label{stacks-properties-lemma-check-immersion-covering}
设 $f : \mathcal{X} \to \mathcal{Y}$ 是代数叠态射。
设 $W$ 是代数空间，且 $W \to \mathcal{Y}$
是满、平坦且局部有限表示的态射。下列条件等价：
\begin{enumerate}
\item $f$ 是（开浸入，分别地，闭浸入）；
\item $V = W \times_\mathcal{Y} \mathcal{X}$ 是代数空间，
且 $V \to W$ 是（开浸入，分别地，闭浸入）。
\end{enumerate}
\end{lemma}

\begin{proof}
这由第 \ref{stacks-properties-section-properties-morphisms} 节的一般讨论推出，
特别是引理 \ref{stacks-properties-lemma-check-representable-covering}
和 \ref{stacks-properties-lemma-check-property-covering}。
\end{proof}

\begin{lemma}
\label{stacks-properties-lemma-immersion-monomorphism}
浸入是单态射。
\end{lemma}

\begin{proof}
见《空间的态射》引理
\ref{spaces-morphisms-lemma-immersions-monomorphisms}。
\end{proof}

\begin{lemma}
\label{stacks-properties-lemma-subspace-induced-topology}
若 $f : \mathcal{X} \to \mathcal{Y}$ 是浸入，则
$|f| : |\mathcal{X}| \to |\mathcal{Y}|$ 是到某个局部闭子集上的同胚。
若 $f$ 是闭浸入，分别地，开浸入，则 $|f|$ 是闭映射，分别地，开映射。
\end{lemma}

\begin{proof}
略。
\end{proof}

\noindent
下面两个引理说明如何利用表示理解浸入。

\begin{lemma}
\label{stacks-properties-lemma-immersion-into-presentation}
设 $(U, R, s, t, c)$ 是代数空间中的光滑群胚，
且 $i : \mathcal{Z} \to [U/R]$ 是浸入。
则存在 $R$-不变的局部闭子空间 $Z \subset U$ 和表示
$[Z/R_Z] \to \mathcal{Z}$，其中 $R_Z$ 是 $R$ 在 $Z$ 上的限制，
使得
$$
\xymatrix{
[Z/R_Z] \ar[dr] \ar[rr] & & \mathcal{Z} \ar[ld]^i \\
& [U/R]
}
$$
为 $2$-交换图表。若 $i$ 是闭浸入（分别地，开浸入），
则 $Z$ 是 $U$ 的闭子空间（分别地，开子空间）。
\end{lemma}

\begin{proof}
由引理 \ref{stacks-properties-lemma-representable-in-terms-presentations}，
得到交换图表
$$
\xymatrix{
[U'/R'] \ar[dr] \ar[rr] & & \mathcal{Z} \ar[ld] \\
& [U/R]
}
$$
其中 $U' = \mathcal{Z} \times_{[U/R]} U$ 且
$R' = \mathcal{Z} \times_{[U/R]} R$。
由于 $\mathcal{Z} \to [U/R]$ 是浸入，$U' \to U$
是代数空间的浸入。令 $Z \subset U$ 为如下局部闭子空间：
$U' \to U$ 经 $Z$ 分解，并诱导同构 $U' \to Z$。
由 $R'$ 的构造显然有
$R' = U' \times_{U, t} R = R \times_{s, U} U'$.
这蕴含 $Z \cong U'$ 是 $R$-不变的，并且 $R' \to R$ 的像将
$R'$ 与 $R$ 在 $Z$ 上的限制
$R_Z = s^{-1}(Z) = t^{-1}(Z)$ 等同。故引理成立。
\end{proof}

\begin{lemma}
\label{stacks-properties-lemma-immersion-presentation}
设 $(U, R, s, t, c)$ 是代数空间中的光滑群胚。
令 $\mathcal{X} = [U/R]$ 为相应的代数叠；见《代数叠》定理
\ref{algebraic-theorem-smooth-groupoid-gives-algebraic-stack}。
设 $Z \subset U$ 是 $R$-不变的局部闭子空间。则
$$
[Z/R_Z] \longrightarrow [U/R]
$$
是代数叠的浸入，其中 $R_Z$ 是 $R$ 在 $Z$ 上的限制。
若 $Z \subset U$ 是开子空间（分别地，闭子空间），则该态射是
代数叠的开浸入（分别地，闭浸入）。
\end{lemma}

\begin{proof}
回忆，由《空间中的群胚》定义
\ref{spaces-groupoids-definition-invariant-open}
（另见该定义之后的讨论），作为 $R$ 的局部闭子空间，有
$R_Z = s^{-1}(Z) = t^{-1}(Z)$。因此，两个态射 $R_Z \to Z$
作为 $s$ 和 $t$ 的基变换都是光滑的。故
$(Z, R_Z, s|_{R_Z}, t|_{R_Z}, c|_{R_Z \times_{s, Z, t} R_Z})$ 是
代数空间中的光滑群胚，并且可见 $[Z/R_Z]$ 是代数叠；
见《代数叠》定理
\ref{algebraic-theorem-smooth-groupoid-gives-algebraic-stack}。
《空间中的群胚》引理
\ref{spaces-groupoids-lemma-criterion-fibre-product}
的假设全部满足，因而得到 $2$-纤维方块
$$
\xymatrix{
Z \ar[d] \ar[r] & [Z/R_Z] \ar[d] \\
U \ar[r] & [U/R]
}
$$
由此和引理 \ref{stacks-properties-lemma-check-representable-covering} 可知，
$[Z/R_Z] \to [U/R]$ 可由代数空间表示；进而由引理
\ref{stacks-properties-lemma-check-property-covering} 可知，右侧竖直箭头是浸入
（分别地，闭浸入；分别地，开浸入），当且仅当左侧竖直箭头如此。
\end{proof}

\noindent
可以如下定义开、闭及局部闭子叠。

\begin{definition}
\label{stacks-properties-definition-substacks}
设 $\mathcal{X}$ 是代数叠。
\begin{enumerate}
\item $\mathcal{X}$ 的一个{\it 开子叠}，是严格满子范畴
$\mathcal{X}' \subset \mathcal{X}$，使得 $\mathcal{X}'$ 是代数叠，
且 $\mathcal{X}' \to \mathcal{X}$ 是开浸入。
\item $\mathcal{X}$ 的一个{\it 闭子叠}，是严格满子范畴
$\mathcal{X}' \subset \mathcal{X}$，使得 $\mathcal{X}'$ 是代数叠，
且 $\mathcal{X}' \to \mathcal{X}$ 是闭浸入。
\item $\mathcal{X}$ 的一个{\it 局部闭子叠}，是严格满子范畴
$\mathcal{X}' \subset \mathcal{X}$，使得 $\mathcal{X}'$ 是代数叠，
且 $\mathcal{X}' \to \mathcal{X}$ 是浸入。
\end{enumerate}
\end{definition}

\noindent
使用该定义时应当谨慎。事实上，若
$f : \mathcal{X} \to \mathcal{Y}$ 是代数叠的等价，
且 $\mathcal{X}' \subset \mathcal{X}$ 是开子叠，
则子范畴 $f(\mathcal{X}')$ 未必是 $\mathcal{Y}$ 的开子叠。
问题在于它可能不是{\it 严格}满子范畴；但这也是唯一的问题。
下面给出形式化陈述。

\begin{lemma}
\label{stacks-properties-lemma-substack-image}
对任一浸入 $i : \mathcal{Z} \to \mathcal{X}$，存在唯一的局部闭子叠
$\mathcal{X}' \subset \mathcal{X}$，使得 $i$ 分解为等价
$i' : \mathcal{Z} \to \mathcal{X}'$ 与包含态射
$\mathcal{X}' \to \mathcal{X}$ 的复合。
若 $i$ 是闭浸入（分别地，开浸入），则 $\mathcal{X}'$
是 $\mathcal{X}$ 的闭子叠（分别地，开子叠）。
\end{lemma}

\begin{proof}
略。
\end{proof}

\begin{lemma}
\label{stacks-properties-lemma-substacks-presentation}
设 $[U/R] \to \mathcal{X}$ 是代数叠的一个表示。存在典范双射
$$
\mathcal{X}\text{ 的局部闭子叠 }\mathcal{Z}
\longrightarrow
U\text{ 的 }R\text{-不变局部闭子空间 }Z
$$
将 $\mathcal{Z}$ 映为 $U \times_\mathcal{X} \mathcal{Z}$。
此外，代数叠态射 $f : \mathcal{Y} \to \mathcal{X}$
经 $\mathcal{Z}$ 分解，当且仅当
$\mathcal{Y} \times_\mathcal{X} U \to U$ 经 $Z$ 分解。
闭子叠和开子叠也有相同结论。
\end{lemma}

\begin{proof}
由引理 \ref{stacks-properties-lemma-immersion-into-presentation} 和
\ref{stacks-properties-lemma-immersion-presentation}，该映射是双射。
若 $\mathcal{Y} \to \mathcal{X}$ 经 $\mathcal{Z}$ 分解，
则其基变换 $\mathcal{Y} \times_\mathcal{X} U \to U$
当然经 $Z$ 分解。反之，设态射 $\mathcal{Y} \to \mathcal{X}$
满足 $\mathcal{Y} \times_\mathcal{X} U \to U$ 经 $Z$ 分解。
我们将证明：对每个概形 $T$ 和每个态射 $T \to \mathcal{Y}$，
若该态射由 $\mathcal{Y}$ 在 $T$ 上纤维范畴中的对象 $y$ 给出，
则 $y$ 实际上属于 $\mathcal{Z}$ 在 $T$ 上的纤维范畴。
事实上，纤维积 $T \times_\mathcal{X} U$ 是代数空间，
且 $T \times_\mathcal{X} U \to T$ 是满的光滑态射。
故存在 fppf 覆盖 $\{T_i \to T\}$，使得对所有 $i$，
$T_i \to T$ 都经 $T \times_\mathcal{X} U \to T$ 分解。
于是 $T_i \to \mathcal{X}$ 经
$\mathcal{Y} \times_\mathcal{X} U$，进而经 $Z \subset U$ 分解。
因此 $y|_{T_i}$ 是 $\mathcal{Z}$ 的对象
（因为 $Z$ 是 $U$ 与 $\mathcal{Z}$ 在 $\mathcal{X}$ 上的纤维积）。
由于 $\mathcal{Z}$ 是严格满子叠，可知 $y$ 是
$\mathcal{Z}$ 的对象，正如所需。
\end{proof}

\begin{lemma}
\label{stacks-properties-lemma-open-substacks}
设 $\mathcal{X}$ 是代数叠。规则
$\mathcal{U} \mapsto |\mathcal{U}|$ 在 $\mathcal{X}$ 的开子叠与
$|\mathcal{X}|$ 的开子集之间定义一个保持包含关系的双射。
\end{lemma}

\begin{proof}
选取表示 $[U/R] \to \mathcal{X}$；见《代数叠》引理
\ref{algebraic-lemma-stack-presentation}。
由引理 \ref{stacks-properties-lemma-substacks-presentation}，
开子叠对应于 $U$ 的 $R$-不变开子概形。
另一方面，引理 \ref{stacks-properties-lemma-points-presentation} 和
\ref{stacks-properties-lemma-topology-points} 保证后者与
$|\mathcal{X}|$ 的开子集双射对应。
\end{proof}

\begin{lemma}
\label{stacks-properties-lemma-open-image-substack}
设 $\mathcal X$ 是代数叠，$U$ 是代数空间，并且
$U \to \mathcal X$ 是满的光滑态射。对任一开浸入
$V \hookrightarrow U$，存在代数叠 $\mathcal Y$、开浸入
$\mathcal Y \to \mathcal X$，以及满的光滑态射
$V \to \mathcal Y$。
\end{lemma}

\begin{proof}
定义群胚纤维化范畴 $\mathcal Y$：对
$(\Sch/S)_{fppf}$ 的对象 $T$，令其上纤维范畴
$\mathcal{Y}_T$ 为 $\mathcal{X}_T$ 的满子范畴，
其对象是所有满足投影态射
$V \times_{\mathcal X, y} T \to T$ 为满射的
$y \in \Ob(\mathcal{X}_T)$。
现对任一态射 $x : T \to \mathcal X$，$2$-纤维积
$T \times_{x, \mathcal X} \mathcal Y$ 在 $T'$ 上的纤维范畴
由三元组 $(f : T' \to T, y \in \mathcal{X}_{T'}, f^*x \simeq y)$
组成，并要求 $V \times_{\mathcal X, y} T' \to T'$ 为满射。
注意，$T \times_{x, \mathcal X} \mathcal Y$ 在集合胚中纤维化，
因为 $\mathcal Y \to \mathcal X$ 忠实
（见《叠》引理
\ref{stacks-lemma-2-fibre-product-gives-stack-in-setoids}）。
同构 $f^*x \simeq y$ 给出图表
$$
\xymatrix{
V \times_{\mathcal X, y} T' \ar[d] \ar[r] &
V \times_{\mathcal X, x} T \ar[r] \ar[d] &
V \ar[d] \\
T' \ar[r]^f &
T \ar[r]^x &
\mathcal X
}
$$
其中两个方块都是笛卡儿的。态射
$V \times_{\mathcal X, x} T \to T$ 作为基变换是光滑的，因而是开的。
令 $T_0 \subset T$ 为其像。由笛卡儿方块可得，
$V \times_{\mathcal X, y} T' \to T'$ 为满射，当且仅当
$f$ 落在 $T_0$ 中。因此，$T \times_{x, \mathcal X} \mathcal Y$
可由 $T_0$ 表示，所以包含 $\mathcal Y \to \mathcal X$ 是开浸入。
由《代数叠》引理
\ref{algebraic-lemma-open-fibred-category-is-algebraic}，
可知 $\mathcal{Y}$ 是代数叠。
最后，将态射 $V \to \mathcal X$ 记作 $g$，则
$V \times_{\mathcal X} V \to V$ 是满射（对角态射给出一个截面）。
故 $g$ 属于 $\mathcal{Y}_V \to \mathcal{X}_V$ 的像；亦即，
得到态射 $g' : V \to \mathcal{Y}$，它嵌入交换图表
$$
\xymatrix{
V \ar[r] \ar[d]^{g'} & U \ar[d] \\
\mathcal{Y} \ar[r] & \mathcal{X}
}
$$
由于 $V \times_{g, \mathcal X} \mathcal Y \to V$ 是单态射，
而 $(1, g')$ 给出一个截面，它实际上是同构。
因此，$g' : V \to \mathcal Y$ 是光滑态射，因为它是光滑态射
$g : V \to \mathcal{X}$ 的基变换。由 $\mathcal{Y}$ 的构造，
它还是满射；引理得证。
\end{proof}

\begin{lemma}
\label{stacks-properties-lemma-union-open-substacks}
设 $\mathcal X$ 是代数叠，且
$\mathcal{X}_i \subset \mathcal X$ 是由 $i \in I$ 标号的一族开子叠。
则存在一个开子叠，记作
$\bigcup_{i\in I} \mathcal{X}_i \subset \mathcal X$，
使得 $\mathcal{X}_i$ 是覆盖它的开子叠。
\end{lemma}

\begin{proof}
定义纤维化子范畴
$\mathcal{X}' = \bigcup_{i \in I} \mathcal{X}_i$：
对 $(\Sch/S)_{fppf}$ 的对象 $T$，令其上纤维范畴为
$\mathcal{X}_T$ 的满子范畴，其对象是所有满足态射
$\coprod_{i \in I} (\mathcal{X}_i \times_{\mathcal X} T) \to T$
为满射的 $x \in \Ob(\mathcal{X}_T)$。
令 $x_i \in \Ob((\mathcal{X}_i)_T)$。则 $(x_i, 1)$ 给出
$\mathcal{X}_i \times_{\mathcal X} T \to T$ 的截面，故有一个同构。
因此 $\mathcal{X}_i \subset \mathcal{X}'$ 是满子范畴。
现取 $x \in \Ob(\mathcal{X}_T)$。则
$\mathcal{X}_i \times_{\mathcal X} T$ 可由开子概形
$T_i \subset T$ 表示。$2$-纤维积
$\mathcal{X}' \times_{\mathcal X} T$ 在 $T'$ 上的纤维
由 $(y \in \mathcal{X}_{T'}, f : T' \to T, f^*x \simeq y)$ 组成，
并要求 $\coprod (\mathcal{X}_i \times_{\mathcal X, y} T') \to T'$
为满射。同构 $f^*x \simeq y$ 诱导同构
$\mathcal{X}_i \times_{\mathcal X, y} T' \simeq T_i \times_T T'$。
于是 $T_i \times_T T'$ 覆盖 $T'$，当且仅当
$f$ 落在 $\bigcup T_i$ 中。因此有图表
$$
\xymatrix{
T_i \ar[r] \ar[d] &
\bigcup T_i \ar[r] \ar[d] &
T \ar[d] \\
\mathcal{X}_i \ar[r] &
\mathcal{X}' \ar[r] &
\mathcal{X}
}
$$
其中两个方块均为笛卡儿的。由《代数叠》引理
\ref{algebraic-lemma-open-fibred-category-is-algebraic}，
可知 $\mathcal{X'} \subset \mathcal{X}$ 是代数的，且为开子叠。
由上述笛卡儿方块还显然可知，态射
$\coprod_{i \in I} \mathcal{X}_i \to \mathcal{X}'$，从而完成引理的证明。
\end{proof}

\begin{lemma}
\label{stacks-properties-lemma-quasi-compact-finite-subcover}
设 $\mathcal X$ 是代数叠，且 $\mathcal X' \subset \mathcal X$
是拟紧开子叠。设有由 $i \in I$ 标号的一族开子叠
$\mathcal{X}_i \subset \mathcal X$，满足
$\mathcal{X}' \subset \bigcup_{i \in I} \mathcal{X}_i$，
其中并按引理 \ref{stacks-properties-lemma-union-open-substacks} 定义。
则存在有限子集 $I' \subset I$，使得
$\mathcal{X}' \subset \bigcup_{i \in I'} \mathcal{X}_i$。
\end{lemma}

\begin{proof}
由于 $\mathcal X$ 是代数的，存在概形 $U$ 和满的光滑态射
$U \to \mathcal X$。令 $U_i \subset U$ 为表示
$\mathcal{X}_i \times_{\mathcal X} U$ 的开子概形，
并令 $U' \subset U$ 为表示
$\mathcal{X}' \times_{\mathcal X} U$ 的开子概形。
由假设，$U'\subset \bigcup_{i\in I} U_i$。
由引理 \ref{stacks-properties-lemma-quasi-compact-stack} 的证明，存在拟紧开子集
$V \subset U'$，使得 $V \to \mathcal{X}'$ 是满的光滑态射。
因此存在有限子集 $I' \subset I$，使得
$V \subset \bigcup_{i \in I'} U_i$。
我们断言 $\mathcal{X}' \subset \bigcup_{i \in I'} \mathcal{X}_i$。
对 $T \in \Ob((\Sch/S)_{fppf})$，取
$x \in \Ob(\mathcal{X}'_T)$。
由于 $\mathcal{X}' \to \mathcal{X}$ 是单态射，有笛卡儿方块
$$
\xymatrix{
V \times_\mathcal{X} T \ar[r] \ar[d] &
T \ar[d]^x \ar@{=}[r] &
T \ar[d]^x \\
V \ar[r] &
\mathcal{X}' \ar[r] &
\mathcal X
}
$$
由基变换，$V \times_{\mathcal X} T \to T$ 是满射。因此
$\bigcup_{i \in I'} U_i \times_{\mathcal X} T \to T$ 也是满射。
令 $T_i \subset T$ 为表示
$\mathcal{X}_i \times_{\mathcal X} T$ 的开子概形。
由形式论证，得到笛卡儿方块
$$
\xymatrix{
U_i \times_{\mathcal{X}_i} T_i \ar[r] \ar[d] &
U \times_{\mathcal X} T \ar[d] \\
T_i \ar[r] & T
}
$$
其中竖直箭头由基变换都是满射。由于
$U_i \times_{\mathcal{X}_i} T_i \simeq U_i \times_{\mathcal X} T$，
可得 $\bigcup_{i \in I'} T_i = T$。因此，由并的定义，
$x$ 是 $(\bigcup_{i\in I'} \mathcal{X}_i)_T$ 的对象。
注意，包含
$\mathcal{X}' \subset \bigcup_{i \in I'} \mathcal{X}_i$
自动是开子叠。
\end{proof}

\begin{lemma}
\label{stacks-properties-lemma-zariski-open-cover-stack-is-space}
设 $\mathcal X$ 是代数叠，且 $\mathcal{X}_i$（$i \in I$）
是一族 $\mathcal{X}$ 的开子叠。假设
\begin{enumerate}
\item $\mathcal{X} = \bigcup_{i \in I} \mathcal{X}_i$；
\item 每个 $\mathcal{X}_i$ 都是代数空间。
\end{enumerate}
则 $\mathcal{X}$ 是代数空间。
\end{lemma}

\begin{proof}
将《叠》引理 \ref{stacks-lemma-stack-in-setoids-descent}
应用于态射 $\coprod_{i \in I} \mathcal{X}_i \to \mathcal{X}$
和态射 $\text{id} : \mathcal{X} \to \mathcal{X}$，
可见 $\mathcal{X}$ 是集合胚中的叠。
因此 $\mathcal{X}$ 是代数空间；见《代数叠》命题
\ref{algebraic-proposition-algebraic-stack-no-automorphisms}。
\end{proof}

\begin{lemma}
\label{stacks-properties-lemma-zariski-open-cover-stack-is-scheme}
设 $\mathcal X$ 是代数叠，且 $\mathcal{X}_i$（$i \in I$）
是一族 $\mathcal{X}$ 的开子叠。假设
\begin{enumerate}
\item $\mathcal{X} = \bigcup_{i \in I} \mathcal{X}_i$；
\item 每个 $\mathcal{X}_i$ 都是概形。
\end{enumerate}
则 $\mathcal{X}$ 是概形。
\end{lemma}

\begin{proof}
由引理 \ref{stacks-properties-lemma-zariski-open-cover-stack-is-space}，
$\mathcal{X}$ 是代数空间。每个代数空间都有一个为概形的最大开子空间；
见《空间的性质》引理
\ref{spaces-properties-lemma-subscheme}。故 $\mathcal{X}$ 是概形。
\end{proof}


\noindent
下一个引理类似于《群胚进阶》引理
\ref{more-groupoids-lemma-local-source}。

\begin{lemma}
\label{stacks-properties-lemma-local-source}
设 $\mathcal{P}, \mathcal{Q}, \mathcal{R}$ 是代数空间态射的性质。假设
\begin{enumerate}
\item $\mathcal{P}, \mathcal{Q}, \mathcal{R}$ 在靶上 fppf 局部，
且在任意基变换下保持；
\item $\text{smooth} \Rightarrow \mathcal{R}$,
\item 对任一具有性质 $\mathcal{Q}$ 的态射 $f : X \to Y$，
存在最大开子空间 $W(\mathcal{P}, f) \subset X$，使得
$f|_{W(\mathcal{P}, f)}$ 具有性质 $\mathcal{P}$；
\item 对任一具有性质 $\mathcal{Q}$ 的态射 $f : X \to Y$，
以及任一具有性质 $\mathcal{R}$ 的态射 $Y' \to Y$，都有
$Y' \times_Y W(\mathcal{P}, f) = W(\mathcal{P}, f')$，
其中 $f' : X_{Y'} \to Y'$ 是 $f$ 的基变换。
\end{enumerate}
设 $f : \mathcal{X} \to \mathcal{Y}$ 是可由代数空间表示的代数叠态射，
并假设 $f$ 具有性质 $\mathcal{Q}$。则
\begin{enumerate}
\item[(A)] 存在最大开子叠 $\mathcal{X}' \subset \mathcal{X}$，
使得 $f|_{\mathcal{X}'}$ 具有性质 $\mathcal{P}$；
\item[(B)] 若 $\mathcal{Z} \to \mathcal{Y}$
是可由代数空间表示且具有性质 $\mathcal{R}$ 的代数叠态射，
则 $\mathcal{Z} \times_\mathcal{Y} \mathcal{X}'$
是 $\mathcal{Z} \times_\mathcal{Y} \mathcal{X}$ 的最大开子叠，
在它上面，基变换 $\text{id}_\mathcal{Z} \times f$
具有性质 $\mathcal{P}$。
\end{enumerate}
\end{lemma}

\begin{proof}
选取概形 $V$ 以及满的光滑态射 $V \to \mathcal{Y}$。
令 $U = V \times_\mathcal{Y} \mathcal{X}$，并令
$f' : U \to V$ 为 $f$ 的基变换。代数空间态射
$f' : U \to V$ 具有性质 $\mathcal{Q}$。
因此，由假设 (3) 得到开子空间 $W(\mathcal{P}, f') \subset U$。
注意
$U \times_\mathcal{X} U =
(V \times_\mathcal{Y} V) \times_\mathcal{Y} \mathcal{X}$
，故态射
$f'' : U \times_\mathcal{X} U \to V \times_\mathcal{Y} V$
是 $f$ 沿任一投影 $V \times_\mathcal{Y} V \to V$ 的基变换。
由 $V$ 的选择，这些投影光滑，故由 (2) 具有性质 $\mathcal{R}$。
因此，由 (4)，$W(\mathcal{P}, f')$ 在两个投影
$\text{pr}_i : U \times_\mathcal{X} U \to U$ 下的逆像相同。
换言之，$W(\mathcal{P}, f')$ 是 $U$ 的 $R$-不变子空间
（其中 $R = U \times_\mathcal{X} U$）。
令 $\mathcal{X}'$ 为通过引理
\ref{stacks-properties-lemma-immersion-into-presentation} 与
$W(\mathcal{P}, f)$ 对应的 $\mathcal{X}$ 的开子叠。
按构造，
$W(\mathcal{P}, f') = \mathcal{X}' \times_\mathcal{Y} V$，
故由引理 \ref{stacks-properties-lemma-check-property-covering}，
$f|_{\mathcal{X}'}$ 具有性质 $\mathcal{P}$。
此外，$\mathcal{X}'$ 是使得 $f|_{\mathcal{X}'}$ 具有
$\mathcal{P}$ 的最大开子叠，因为 $W(\mathcal{P}, f)$
也具有同样的极大性。这证明了 (A)。

\medskip\noindent
最后，若 $\mathcal{Z} \to \mathcal{Y}$ 是可由代数空间表示且
具有性质 $\mathcal{R}$ 的代数叠态射，则令
$T = V \times_\mathcal{Y} \mathcal{Z}$。
可见 $T \to V$ 是具有性质 $\mathcal{R}$ 的代数空间态射。
令 $f'_T : T \times_V U \to T$ 为 $f'$ 的基变换。
再次由 (4)，$W(\mathcal{P}, f'_T)$ 是
$W(\mathcal{P}, f)$ 在 $T \times_V U$ 中的逆像。
这蕴含 (B)；略去一些细节。
\end{proof}

\begin{remark}
\label{stacks-properties-remark-local-source-warning}
警告：使用引理 \ref{stacks-properties-lemma-local-source} 时应当谨慎。例如，它适用于
$\mathcal{P}=$``平坦''、$\mathcal{Q}=$``空性质''、
$\mathcal{R}=$``平坦且局部有限表示''的情形。
但给定代数空间态射 $f : X \to Y$，使得 $f|_W$ 平坦的最大开子空间
$W \subset X$ {\it 并不是} $f$ 平坦的点集！
\end{remark}

\begin{remark}
\label{stacks-properties-remark-local-source-apply}
尽管有注 \ref{stacks-properties-remark-local-source-warning} 中的警告，
在某些情形下使用引理 \ref{stacks-properties-lemma-local-source} 不会引起歧义。
下面列出这些情形。每种情形中，我们略去假设 (1) 和 (2) 的验证，
并给出蕴含 (3) 和 (4) 的参考文献。具体如下：
\begin{enumerate}
\item
\label{stacks-properties-item-rel-dim-leq-d}
$\mathcal{Q} = $``局部有限型''，$\mathcal{R} = \emptyset$，
$\mathcal{P} =$``相对维数 $\leq d$''。
见《空间的态射》定义
\ref{spaces-morphisms-definition-relative-dimension}
以及《空间的态射》引理
\ref{spaces-morphisms-lemma-openness-bounded-dimension-fibres} 和
\ref{spaces-morphisms-lemma-dimension-fibre-after-base-change}。
\item
\label{stacks-properties-item-loc-quasi-finite}
$\mathcal{Q} =$``局部有限型''，$\mathcal{R} = \emptyset$，
$\mathcal{P} =$``局部拟有限''。
这是上一项中 $d = 0$ 的情形；见《空间的态射》引理
\ref{spaces-morphisms-lemma-locally-quasi-finite-rel-dimension-0}.
另一方面，性质 (3) 和 (4) 在《空间的态射》引理
\ref{spaces-morphisms-lemma-locally-finite-type-quasi-finite-part}
中明确给出。
\item
\label{stacks-properties-item-unramified}
$\mathcal{Q} = $``局部有限型''，$\mathcal{R} = \emptyset$，
$\mathcal{P} =$``非分歧''。这就是《空间的态射》引理
\ref{spaces-morphisms-lemma-where-unramified}。
\item
\label{stacks-properties-item-flat}
$\mathcal{Q} =$``局部有限表示''，
$\mathcal{R} =$``平坦且局部有限表示''，
$\mathcal{P} =$``平坦''。见《空间态射进阶》定理
\ref{spaces-more-morphisms-theorem-openness-flatness} 和引理
\ref{spaces-more-morphisms-lemma-flat-locus-base-change}。
注意，此处 $W(\mathcal{P}, f)$ 总是恰好等于态射 $f$ 平坦的点集，
因为只在 $f$ 具有性质 $\mathcal{Q}$ 时才考虑这个开集
（见同处引文）。
\item
\label{stacks-properties-item-etale}
$\mathcal{Q} =$``局部有限表示''，
$\mathcal{R} =$``平坦且局部有限表示''，
$\mathcal{P}=$``\'etale''。这由合并
(\ref{stacks-properties-item-unramified}) 和 (\ref{stacks-properties-item-flat}) 得出，
因为平坦、局部有限表示的非分歧态射是 \'etale 态射；
见《空间的态射》引理
\ref{spaces-morphisms-lemma-unramified-flat-lfp-etale}。
\item 按需在此添加更多情形（与《群胚进阶》注
\ref{more-groupoids-remark-local-source-apply} 中更长的列表比较）。
\end{enumerate}
\end{remark}











\section{既约代数叠}
\label{stacks-properties-section-reduced}

\noindent
我们已经在第 \ref{stacks-properties-section-types-properties} 节定义了既约代数叠。

\begin{lemma}
\label{stacks-properties-lemma-reduced-closed-substack}
设 $\mathcal{X}$ 是代数叠，且 $T \subset |\mathcal{X}|$ 是闭子集。
存在唯一的闭子叠 $\mathcal{Z} \subset \mathcal{X}$，满足：
(a) $|\mathcal{Z}| = T$；(b) $\mathcal{Z}$ 既约。
\end{lemma}

\begin{proof}
设 $U \to \mathcal{X}$ 是满的光滑态射，其中 $U$ 是代数空间。
令 $R = U \times_\mathcal{X} U$，于是存在表示
$[U/R] \to \mathcal{X}$；见《代数叠》引理
\ref{algebraic-lemma-stack-presentation}。
照例将两个光滑投影态射记作 $s, t : R \to U$。
由引理 \ref{stacks-properties-lemma-points-presentation} 可知，$T$ 对应于闭子集
$T' \subset |U|$，并且
$|s|^{-1}(T') = |t|^{-1}(T')$。
令 $Z \subset U$ 为 $T'$ 上的既约诱导代数空间结构；
见《空间的性质》定义
\ref{spaces-properties-definition-reduced-induced-space}。
纤维积 $Z \times_{U, t} R$ 和 $R \times_{s, U} Z$
是 $R$ 的闭子空间
（见《空间》引理 \ref{spaces-lemma-base-change-immersions}）。
由《空间的态射》引理
\ref{spaces-morphisms-lemma-base-change-smooth}，投影
$Z \times_{U, t} R \to Z$ 和 $R \times_{s, U} Z \to Z$
都是光滑的。因此，由于 $Z$ 既约，注
\ref{stacks-properties-remark-list-properties-local-smooth-topology} 蕴含
$Z \times_{U, t} R$ 和 $R \times_{s, U} Z$ 都既约。
由于
$$
|Z \times_{U, t} R| = |t|^{-1}(T') = |s|^{-1}(T') = R \times_{s, U} Z
$$
由《空间的性质》引理
\ref{spaces-properties-lemma-reduced-closed-subspace} 中的唯一性，
可得 $Z \times_{U, t} R = R \times_{s, U} Z$。
故 $Z$ 是 $U$ 的 $R$-不变闭子空间。
通过引理 \ref{stacks-properties-lemma-substacks-presentation} 的对应，
得到闭子叠 $\mathcal{Z} \subset \mathcal{X}$，满足
$Z = \mathcal{Z} \times_\mathcal{X} U$。于是
$[Z/R_Z] \to \mathcal{Z}$ 是一个表示
（引理 \ref{stacks-properties-lemma-immersion-into-presentation}）。
因此 $|\mathcal{Z}| = |Z|/|R_Z| = |T'|/\sim$
就是给定闭子集 $T$。略去唯一性的证明。
\end{proof}

\begin{lemma}
\label{stacks-properties-lemma-reduced-stack-determined-by-points}
设 $\mathcal{X}$ 是代数叠。若
$\mathcal{X}' \subset \mathcal{X}$ 是闭子叠，
$\mathcal{X}$ 既约且
$|\mathcal{X}'| = |\mathcal{X}|$，则
$\mathcal{X}' = \mathcal{X}$。
\end{lemma}

\begin{proof}
选取表示 $[U/R] \to \mathcal{X}$，其中 $U$ 是概形。
由于 $\mathcal{X}$ 既约，由既约代数叠的定义，$U$ 既约。
由引理 \ref{stacks-properties-lemma-substacks-presentation}，
$\mathcal{X}'$ 对应于 $R$-不变闭子概形 $Z \subset U$。
现在 $|Z| \subset |U|$ 是 $|\mathcal{X}'|$ 的逆像，因而
$|Z| = |U|$。故 $Z$ 是 $U$ 的闭子概形，二者的底层点集相同。
由《概形》引理 \ref{schemes-lemma-map-into-reduction}，
映射 $\text{id}_U : U \to U$ 经 $Z \to U$ 分解，故
$Z = U$，亦即 $\mathcal{X}' = \mathcal{X}$。
\end{proof}

\begin{lemma}
\label{stacks-properties-lemma-map-into-reduction}
设 $\mathcal{X}$、$\mathcal{Y}$ 是代数叠。
设 $\mathcal{Z} \subset \mathcal{X}$ 是闭子叠。
假设 $\mathcal{Y}$ 既约。态射
$f : \mathcal{Y} \to \mathcal{X}$ 经 $\mathcal{Z}$ 分解，
当且仅当 $f(|\mathcal{Y}|) \subset |\mathcal{Z}|$。
\end{lemma}

\begin{proof}
假设 $f(|\mathcal{Y}|) \subset |\mathcal{Z}|$。考虑
$\mathcal{Y} \times_\mathcal{X} \mathcal{Z} \to \mathcal{Y}$。
存在等价
$\mathcal{Y} \times_\mathcal{X} \mathcal{Z} \to \mathcal{Y}'$，
其中 $\mathcal{Y}'$ 是 $\mathcal{Y}$ 的闭子叠；
见引理 \ref{stacks-properties-lemma-base-change-immersion} 和
\ref{stacks-properties-lemma-substack-image}。
利用引理 \ref{stacks-properties-lemma-points-cartesian}、
\ref{stacks-properties-lemma-monomorphism-injective-points} 和
\ref{stacks-properties-lemma-immersion-monomorphism}，可见
$|\mathcal{Y}'| = |\mathcal{Y}|$。因此，本引理归结为
引理 \ref{stacks-properties-lemma-reduced-stack-determined-by-points}。
\end{proof}

\begin{definition}
\label{stacks-properties-definition-reduced-induced-stack}
设 $\mathcal{X}$ 是代数叠，且
$Z \subset |\mathcal{X}|$ 是闭子集。
$Z$ 上的一个{\it 代数叠结构}，由 $\mathcal{X}$ 的一个闭子叠
$\mathcal{Z}$ 给出，并要求 $|\mathcal{Z}|$ 等于 $Z$。
$Z$ 上的{\it 既约诱导代数叠结构}，是引理
\ref{stacks-properties-lemma-reduced-closed-substack} 中构造的结构。
$\mathcal{X}$ 的{\it 既约化 $\mathcal{X}_{red}$}，
是 $|\mathcal{X}|$ 上的既约诱导代数叠结构。
\end{definition}

\noindent
事实上，可以用它定义局部闭子集上的既约诱导代数叠结构。

\begin{remark}
\label{stacks-properties-remark-stack-structure-locally-closed-subset}
设 $X$ 是代数叠，且
$T \subset |\mathcal{X}|$ 是局部闭子集。
令 $\partial T$ 为 $T$ 在拓扑空间 $|\mathcal{X}|$ 中的边界。用公式表示为
$$
\partial T = \overline{T} \setminus T.
$$
令 $\mathcal{U} \subset \mathcal{X}$ 为 $X$ 的开子叠，满足
$|\mathcal{U}| = |\mathcal{X}| \setminus \partial T$；
见引理 \ref{stacks-properties-lemma-open-substacks}。
令 $\mathcal{Z}$ 为 $\mathcal{U}$ 的既约闭子叠，满足
$|\mathcal{Z}| = T$；它通过取既约诱导闭子空间结构得到，
见定义 \ref{stacks-properties-definition-reduced-induced-stack}。
按构造，$\mathcal{Z} \to \mathcal{U}$ 是代数叠的闭浸入，
而 $\mathcal{U} \to \mathcal{X}$ 是开浸入。
因此，由引理 \ref{stacks-properties-lemma-composition-immersion}，
$\mathcal{Z} \to \mathcal{X}$ 是代数叠的浸入。
注意，$\mathcal{Z}$ 是既约代数叠，并且作为 $|X|$ 的子集，
$|\mathcal{Z}| = T$。有时称 $\mathcal{Z}$ 为 $T$ 上的
{\it 既约诱导子叠结构}。
\end{remark}













\section{剩余胚}
\label{stacks-properties-section-residual-gerbe}

\noindent
在 Stacks 项目中，我们希望将代数叠 $\mathcal{X}$ 在点
$x \in |\mathcal{X}|$ 处的{\it 剩余胚}定义为代数叠的单态射
$m_x : \mathcal{Z}_x \to \mathcal{X}$，其中
$\mathcal{Z}_x$ 是仅有一个点的既约代数叠，且该点在 $m_x$ 下映到 $x$。
事实证明，这一概念有许多问题：一般而言，其存在性不明确，唯一性也不明确。
我们通过对代数叠 $\mathcal{Z}_x$ 施加稍强的条件来解决唯一性问题。
下面通过几个关于仅有一个点的既约代数叠的简单引理作更详细的讨论。

\begin{lemma}
\label{stacks-properties-lemma-flat-cover-by-field}
设 $\mathcal{Z}$ 是代数叠，$k$ 是域，并设
$\Spec(k) \to \mathcal{Z}$ 满且平坦。
则对任一域 $k'$，任一态射
$\Spec(k') \to \mathcal{Z}$ 都满且平坦。
\end{lemma}

\begin{proof}
考虑纤维方块
$$
\xymatrix{
T \ar[d] \ar[r] & \Spec(k) \ar[d] \\
\Spec(k') \ar[r] & \mathcal{Z}
}
$$
注意，$T \to \Spec(k')$ 平坦且满，故 $T$ 非空。
另一方面，由于 $k$ 是域，$T \to \Spec(k)$ 平坦。
因此 $T \to \mathcal{Z}$ 平坦且满。
由《空间的态射》引理
\ref{spaces-morphisms-lemma-flat-permanence}
（通过第 \ref{stacks-properties-section-properties-morphisms} 节的讨论），
可知 $\Spec(k') \to \mathcal{Z}$ 平坦。
又因假设 $|\mathcal{Z}|$ 是单点集，它显然是满射。
\end{proof}

\begin{lemma}
\label{stacks-properties-lemma-unique-point}
设 $\mathcal{Z}$ 是代数叠。下列条件等价：
\begin{enumerate}
\item $\mathcal{Z}$ 既约且 $|\mathcal{Z}|$ 是单点集；
\item 存在满的平坦态射 $\Spec(k) \to \mathcal{Z}$，其中 $k$ 是域；
\item 存在局部有限型、满且平坦的态射
$\Spec(k) \to \mathcal{Z}$，其中 $k$ 是域。
\end{enumerate}
\end{lemma}

\begin{proof}
设 (1) 成立。令 $W$ 为概形，并设
$W \to \mathcal{Z}$ 是满的光滑态射。则 $W$ 是既约概形。
令 $\eta \in W$ 为 $W$ 某个不可约分支的泛点。
由于 $W$ 既约，有 $\mathcal{O}_{W, \eta} = \kappa(\eta)$。
因此典范态射 $\eta = \Spec(\kappa(\eta)) \to W$ 平坦。
于是复合 $\eta \to \mathcal{Z}$ 平坦
（见《空间的态射》引理
\ref{spaces-morphisms-lemma-composition-flat}）。
又因 $|\mathcal{Z}|$ 是单点集，该复合也是满射。换言之，(2) 成立。

\medskip\noindent
设 (2) 成立。令 $W$ 为概形，并设
$W \to \mathcal{Z}$ 是满的光滑态射。选取域 $k$
和满的平坦态射 $\Spec(k) \to \mathcal{Z}$。
则 $W \times_\mathcal{Z} \Spec(k)$ 是 $k$ 上光滑的代数空间，
因而正则（见《域上的空间》引理
\ref{spaces-over-fields-lemma-smooth-regular}），特别地既约。
由于 $W \times_\mathcal{Z} \Spec(k) \to W$ 满且平坦，
由《空间上的下降》引理
\ref{spaces-descent-lemma-descend-reduced} 可知 $W$ 既约。
换言之，(1) 成立。

\medskip\noindent
显然，(3) 蕴含 (2)。最后，设 (2) 成立。
选取非空仿射概形 $W$ 和光滑态射 $W \to \mathcal{Z}$。
选取闭点 $w \in W$，并令 $k = \kappa(w)$。复合
$$
\Spec(k) \xrightarrow{w} W \longrightarrow \mathcal{Z}
$$
由《空间的态射》引理
\ref{spaces-morphisms-lemma-composition-finite-type} 和
\ref{spaces-morphisms-lemma-smooth-locally-finite-type}.
为局部有限型。由引理 \ref{stacks-properties-lemma-flat-cover-by-field}，
它还平坦且满。因此 (3) 成立。
\end{proof}

\noindent
下一个引理从仅有一个点的代数叠中，挑出一个比上一引理略好的类。

\begin{lemma}
\label{stacks-properties-lemma-unique-point-better}
设 $\mathcal{Z}$ 是代数叠。下列条件等价：
\begin{enumerate}
\item $\mathcal{Z}$ 既约、局部 Noether，且
$|\mathcal{Z}|$ 是单点集；
\item 存在局部有限表示、满且平坦的态射
$\Spec(k) \to \mathcal{Z}$，其中 $k$ 是域。
\end{enumerate}
\end{lemma}

\begin{proof}
设 (2) 成立。由引理 \ref{stacks-properties-lemma-unique-point} 可知，
$\mathcal{Z}$ 既约且 $|\mathcal{Z}|$ 是单点集。
令 $W$ 为概形，并设 $W \to \mathcal{Z}$ 是满的光滑态射。
选取域 $k$ 以及局部有限表示、满且平坦的态射
$\Spec(k) \to \mathcal{Z}$。
则 $W \times_\mathcal{Z} \Spec(k)$ 是 $k$ 上光滑的代数空间，
因而局部 Noether（见《空间的态射》引理
\ref{spaces-morphisms-lemma-locally-finite-type-locally-noetherian}）。
由于 $W \times_\mathcal{Z} \Spec(k) \to W$
平坦、满且局部有限表示，
$\{W \times_\mathcal{Z} \Spec(k) \to W\}$ 是 fppf 覆盖。
由《空间上的下降》引理
\ref{spaces-descent-lemma-descend-locally-Noetherian}，
可知 $W$ 局部 Noether。换言之，(1) 成立。

\medskip\noindent
设 (1) 成立。选取非空仿射概形 $W$ 和光滑态射
$W \to \mathcal{Z}$。选取闭点 $w \in W$，并令
$k = \kappa(w)$。由于 $W$ 局部 Noether，态射
$w : \Spec(k) \to W$ 是有限表示的；见《态射》引理
\ref{morphisms-lemma-closed-immersion-finite-presentation}。
因此复合
$$
\Spec(k) \xrightarrow{w} W \longrightarrow \mathcal{Z}
$$
由《空间的态射》引理
\ref{spaces-morphisms-lemma-composition-finite-presentation} 和
\ref{spaces-morphisms-lemma-smooth-locally-finite-presentation}.
为局部有限表示。由引理 \ref{stacks-properties-lemma-flat-cover-by-field}，
它还平坦且满。因此 (2) 成立。
\end{proof}

\begin{lemma}
\label{stacks-properties-lemma-monomorphism-into-point}
设 $\mathcal{Z}' \to \mathcal{Z}$ 是代数叠的单态射。
假设存在域 $k$ 和局部有限表示、满且平坦的态射
$\Spec(k) \to \mathcal{Z}$。则 $\mathcal{Z}'$ 或为空，
或 $\mathcal{Z}' \to \mathcal{Z}$ 是等价。
\end{lemma}

\begin{proof}
可以假设 $\mathcal{Z}'$ 非空。此时纤维积
$T = \mathcal{Z}' \times_\mathcal{Z} \Spec(k)$ 非空；
见引理 \ref{stacks-properties-lemma-points-cartesian}。
现在 $T$ 是代数空间，且投影 $T \to \Spec(k)$ 是单态射。
故 $T = \Spec(k)$；见《空间的态射》引理
\ref{spaces-morphisms-lemma-monomorphism-toward-field}。
于是 $\Spec(k) \to \mathcal{Z}$ 经 $\mathcal{Z}'$ 分解。
设态射 $z : \Spec(k) \to \mathcal{Z}$
由 $\Spec(k)$ 上的对象 $\xi$ 给出。刚才已经看到，
$\xi$ 同构于 $\mathcal{Z}'$ 在 $\Spec(k)$ 上的一个对象 $\xi'$。
由于 $z$ 满、平坦且局部有限表示，
$\mathcal{Z}$ 在任一概形上的每个对象都在 fppf 局部同构于
$\xi$ 的拉回，因而也同构于 $\xi'$ 的拉回。
由群胚叠的对象下降，这蕴含
$\mathcal{Z}' \to \mathcal{Z}$ 本质满
（它也全忠实，见引理 \ref{stacks-properties-lemma-monomorphism}）。结论得证。
\end{proof}

\begin{lemma}
\label{stacks-properties-lemma-improve-unique-point}
设 $\mathcal{Z}$ 是代数叠。假设 $\mathcal{Z}$ 满足引理
\ref{stacks-properties-lemma-unique-point} 的等价条件。
则存在唯一的严格满子范畴
$\mathcal{Z}' \subset \mathcal{Z}$，使得
$\mathcal{Z}'$ 是满足引理
\ref{stacks-properties-lemma-unique-point-better} 之等价条件的代数叠。
包含态射 $\mathcal{Z}' \to \mathcal{Z}$ 是代数叠的单态射。
\end{lemma}

\begin{proof}
最后一个断言由第一个断言和引理 \ref{stacks-properties-lemma-monomorphism} 立即可得。
选取域 $k$ 和满、平坦且局部有限型的态射
$\Spec(k) \to \mathcal{Z}$。
令 $U = \Spec(k)$ 且 $R = U \times_\mathcal{Z} U$。
投影 $s, t : R \to U$ 局部有限型。
由于 $U$ 是域的谱，由《空间的态射》引理
\ref{spaces-morphisms-lemma-noetherian-finite-type-finite-presentation}，
可知 $s, t$ 平坦且局部有限表示。
由《可表示性判据》定理
\ref{criteria-theorem-flat-groupoid-gives-algebraic-stack}，
$\mathcal{Z}' = [U/R]$ 是代数叠。
由《代数叠》引理
\ref{algebraic-lemma-map-space-into-stack}，得到典范态射
$$
f : \mathcal{Z}' \longrightarrow \mathcal{Z}
$$
且它全忠实。因此，该态射可由代数空间表示
（见《代数叠》引理
\ref{algebraic-lemma-characterize-representable-by-algebraic-spaces}），
并且是单态射（见引理 \ref{stacks-properties-lemma-monomorphism}）。
由《可表示性判据》引理
\ref{criteria-lemma-flat-quotient-flat-presentation}，
态射 $U \to \mathcal{Z}'$ 满、平坦且局部有限表示。
故 $\mathcal{Z}'$ 是满足引理
\ref{stacks-properties-lemma-unique-point-better} 之等价条件的代数叠。
由《代数叠》引理 \ref{algebraic-lemma-equivalent}，
可以用 $\mathcal{Z}'$ 在 $\mathcal{Z}$ 中的本质像替代它。
因此，除了 $\mathcal{Z}' \subset \mathcal{Z}$ 的唯一性外，
引理的所有断言均已证明。设
$\mathcal{Z}'' \subset \mathcal{Z}$ 是另一个这样的代数叠。
则投影
$$
\mathcal{Z}'
\longleftarrow
\mathcal{Z}' \times_\mathcal{Z} \mathcal{Z}''
\longrightarrow
\mathcal{Z}''
$$
都是单态射。由引理 \ref{stacks-properties-lemma-points-cartesian}，
中间的代数叠非空。因此，由引理
\ref{stacks-properties-lemma-monomorphism-into-point}，两个投影均为同构，结论得证。
\end{proof}

\begin{example}
\label{stacks-properties-example-distinct}
下面给出一个引理 \ref{stacks-properties-lemma-improve-unique-point}
中所构造态射不是同构的例子。这个例子说明，在定义
\ref{stacks-properties-definition-residual-gerbe} 中要求剩余胚局部 Noether 是必要的。
事实上，该例甚至是代数空间！
令 $\text{Gal}(\overline{\mathbf{Q}}/\mathbf{Q})$ 为
$\mathbf{Q}$ 赋予有限群逆极限拓扑后的绝对 Galois 群。令
$$
U =
\Spec(\overline{\mathbf{Q}})
\times_{\Spec(\mathbf{Q})}
\Spec(\overline{\mathbf{Q}}) =
\text{Gal}(\overline{\mathbf{Q}}/\mathbf{Q}) \times
\Spec(\overline{\mathbf{Q}})
$$
（略去对最后一个等号含义的精确解释）。
令 $G$ 表示赋予离散拓扑的绝对 Galois 群
$\text{Gal}(\overline{\mathbf{Q}}/\mathbf{Q})$，
并将其视为 $\Spec(\overline{\mathbf{Q}})$ 上的常群概形；
见《群胚》例 \ref{groupoids-example-constant-group}。
于是 $G$ 自由且传递地作用在 $U$ 上。
令 $X = U/G$；见《空间》定义 \ref{spaces-definition-quotient}。
则 $X$ 是恰有一个点的非 Noether 既约代数空间。
此外，$X$ 有一个（局部）有限型的点：
$$
x :
\Spec(\overline{\mathbf{Q}}) \longrightarrow
U \longrightarrow X
$$
事实上，$U$ 的每个点都是闭点！
由于 $X$ 是 $\overline{\mathbf{Q}}$ 上的代数空间，
$x$ 是单态射。因此，$x$ 是引理
\ref{stacks-properties-lemma-improve-unique-point} 中构造的态射，但 $x$ 不是同构。
事实上，$\Spec(\overline{\mathbf{Q}}) \to X$
是 $X$ 在 $x$ 处的剩余胚。
\end{example}

\noindent
稍后将会看到，在对代数叠 $\mathcal{X}$ 的温和假设下，
下一个引理的等价条件对每个点 $x \in |\mathcal{X}|$ 都成立
（见《叠的态射》第
\ref{stacks-morphisms-section-existence-residual-gerbes} 节）。

\begin{lemma}
\label{stacks-properties-lemma-residual-gerbe}
设 $\mathcal{X}$ 是代数叠，且 $x \in |\mathcal{X}|$ 是一个点。
下列条件等价：
\begin{enumerate}
\item 存在代数叠 $\mathcal{Z}$ 和单态射
$\mathcal{Z} \to \mathcal{X}$，使得 $|\mathcal{Z}|$ 是单点集，
且 $|\mathcal{Z}|$ 在 $|\mathcal{X}|$ 中的像是 $x$；
\item 存在既约代数叠 $\mathcal{Z}$ 和单态射
$\mathcal{Z} \to \mathcal{X}$，使得 $|\mathcal{Z}|$ 是单点集，
且 $|\mathcal{Z}|$ 在 $|\mathcal{X}|$ 中的像是 $x$；
\item 存在代数叠 $\mathcal{Z}$、单态射
$f : \mathcal{Z} \to \mathcal{X}$，以及满的平坦态射
$z : \Spec(k) \to \mathcal{Z}$，其中 $k$ 是域，并且
$x = f(z)$。
\end{enumerate}
此外，若这些条件成立，则存在唯一的严格满子范畴
$\mathcal{Z}_x \subset \mathcal{X}$，使得
$\mathcal{Z}_x$ 是既约、局部 Noether 的代数叠，
且 $|\mathcal{Z}_x|$ 是单点集，并在映射
$|\mathcal{Z}_x| \to |\mathcal{X}|$ 下映到 $x$。
\end{lemma}

\begin{proof}
若 $\mathcal{Z} \to \mathcal{X}$ 如 (1)，则
$\mathcal{Z}_{red} \to \mathcal{X}$ 如 (2)。
（代数叠既约化的概念见第 \ref{stacks-properties-section-reduced} 节。）
故 (1) 蕴含 (2)。(2) 蕴含 (1) 是直接的。
(2) 与 (3) 的等价性由引理 \ref{stacks-properties-lemma-unique-point} 立即可得。

\medskip\noindent
至此已证明 (1) -- (3) 等价。选取 (2) 中的单态射
$f : \mathcal{Z} \to \mathcal{X}$。
注意，这蕴含 $f$ 全忠实；见引理 \ref{stacks-properties-lemma-monomorphism}。
将函子 $f$ 的本质像记作
$\mathcal{Z}' \subset \mathcal{X}$。则
$f : \mathcal{Z} \to \mathcal{Z}'$ 是等价，故
$\mathcal{Z}'$ 是代数叠；见《代数叠》引理
\ref{algebraic-lemma-equivalent}。
应用引理 \ref{stacks-properties-lemma-improve-unique-point}，
得到引理陈述中的严格满子范畴
$\mathcal{Z}_x \subset \mathcal{Z}'$。
这证明了除唯一性之外的所有断言。

\medskip\noindent
为证明唯一性，设
$\mathcal{Z}_x \subset \mathcal{X}$
和
$\mathcal{Z}'_x \subset \mathcal{X}$
是引理陈述中的两个严格满子范畴。则投影
$$
\mathcal{Z}'_x
\longleftarrow
\mathcal{Z}'_x \times_\mathcal{X} \mathcal{Z}_x
\longrightarrow
\mathcal{Z}_x
$$
都是单态射。由引理 \ref{stacks-properties-lemma-points-cartesian}，
中间的代数叠非空。因此，由引理
\ref{stacks-properties-lemma-monomorphism-into-point}，两个投影均为同构，结论得证。
\end{proof}

\noindent
经过上述说明，现在可以作如下定义。

\begin{definition}
\label{stacks-properties-definition-residual-gerbe}
设 $\mathcal{X}$ 是代数叠，且 $x \in |\mathcal{X}|$。
\begin{enumerate}
\item 若引理 \ref{stacks-properties-lemma-residual-gerbe} 的等价条件
(1)、(2)、(3) 成立，则称{\it $\mathcal{X}$ 在 $x$ 处的剩余胚存在}。
\item 若 $\mathcal{X}$ 在 $x$ 处的剩余胚存在，则
{\it $\mathcal{X}$ 在 $x$ 处的剩余胚}\footnote{这在精神上与
\cite{LM-B} 有冲突，但事实上并无冲突。具体地，他们在第 11 章
为任一拟分离代数叠上的任一点赋予一个胚（不必是代数的），
并称其为剩余胚。我们将在《叠的态射》引理
\ref{stacks-morphisms-lemma-every-point-residual-gerbe}
中看到，在拟分离代数叠上，每个点都有我们意义下的剩余胚，
且该剩余胚与他们的定义等价。关于这一主题的更多信息，见
\cite[Appendix B]{rydh_etale_devissage}。}
是引理 \ref{stacks-properties-lemma-residual-gerbe} 中构造的严格满子范畴
$\mathcal{Z}_x \subset \mathcal{X}$。
\end{enumerate}
\end{definition}

\noindent
特别地，我们知道 $\mathcal{Z}_x$（若存在）是局部 Noether 的
既约代数叠，并且存在一个域和一个满、平坦、局部有限表示的态射
$$
\Spec(k) \longrightarrow \mathcal{Z}_x.
$$
我们将在《叠的态射》引理
\ref{stacks-morphisms-lemma-noetherian-singleton-stack-gerbe}
中看到，$\mathcal{Z}_x$ 是胚。剩余胚的存在性在
《叠的态射》第
\ref{stacks-morphisms-section-existence-residual-gerbes} 节讨论。

\begin{example}
\label{stacks-properties-example-residual-gerbe}
设 $X$ 是概形，且 $x \in X$ 是一个点。则单态射
$x \to X$ 是 $X$ 在 $x$ 处的剩余胚；这里照例将
$x$ 与概形 $x = \Spec(\kappa(x))$ 等同。
若 $X$ 是代数空间且 $x \in |X|$，则 $x$ 处的剩余胚
（称为剩余空间）总是存在；见《适度空间》第
\ref{decent-spaces-section-singleton} 节。
\end{example}

\noindent
由下一个引理，剩余胚若存在，便是正则代数叠。

\begin{lemma}
\label{stacks-properties-lemma-residual-gerbe-regular}
若既约、局部 Noether 的代数叠 $\mathcal{Z}$
满足 $|\mathcal{Z}|$ 是单点集，则它正则。
\end{lemma}

\begin{proof}
设 $W \to \mathcal{Z}$ 是满的光滑态射，其中 $W$ 是概形。
设 $k$ 是域，并设 $\Spec(k) \to \mathcal{Z}$
满、平坦且局部有限表示（见引理 \ref{stacks-properties-lemma-unique-point-better}）。
代数空间 $T = W \times_\mathcal{Z} \Spec(k)$ 在 $k$ 上光滑，
特别地正则；见《域上的空间》引理
\ref{spaces-over-fields-lemma-smooth-regular}。
由于 $T \to W$ 局部有限表示、平坦且满，
由《空间上的下降》引理
\ref{spaces-descent-lemma-descend-regular}，可知 $W$ 正则。
按定义，这意味着 $\mathcal{Z}$ 正则。
\end{proof}

\begin{lemma}
\label{stacks-properties-lemma-residual-gerbe-points}
设 $\mathcal{X}$ 是代数叠，且 $x \in |\mathcal{X}|$。
假设 $\mathcal{X}$ 的剩余胚 $\mathcal{Z}_x$ 存在。
设 $f : \Spec(K) \to \mathcal{X}$ 是 $x$ 的等价类中的态射，
其中 $K$ 是域。则 $f$ 经包含态射
$\mathcal{Z}_x \to \mathcal{X}$ 分解。
\end{lemma}

\begin{proof}
选取域 $k$ 和满、平坦、局部有限表示的态射
$\Spec(k) \to \mathcal{Z}_x$。令
$T = \Spec(K) \times_\mathcal{X} \mathcal{Z}_x$。
由引理 \ref{stacks-properties-lemma-points-cartesian} 可知 $T$ 非空。
由于 $\mathcal{Z}_x \to \mathcal{X}$ 是单态射，
$T \to \Spec(K)$ 是单态射。因此，由《空间的态射》引理
\ref{spaces-morphisms-lemma-monomorphism-toward-field}，
可知 $T = \Spec(K)$，引理得证。
\end{proof}

\begin{lemma}
\label{stacks-properties-lemma-residual-gerbe-unique}
设 $\mathcal{X}$ 是代数叠，且 $x \in |\mathcal{X}|$。
设 $\mathcal{Z}$ 是满足引理
\ref{stacks-properties-lemma-unique-point-better} 之等价条件的代数叠，
并设 $\mathcal{Z} \to \mathcal{X}$ 是单态射，使得
$|\mathcal{Z}| \to |\mathcal{X}|$ 的像是 $x$。
则 $\mathcal{X}$ 在 $x$ 处的剩余胚
$\mathcal{Z}_x$ 存在，并且
$\mathcal{Z} \to \mathcal{X}$ 分解为
$\mathcal{Z} \to \mathcal{Z}_x \to \mathcal{X}$，
其中第一个箭头是等价。
\end{lemma}

\begin{proof}
令 $\mathcal{Z}_x \subset \mathcal{X}$ 为与函子
$\mathcal{Z} \to \mathcal{X}$ 的本质像对应的满子范畴。
则 $\mathcal{Z} \to \mathcal{Z}_x$ 是等价，故
$\mathcal{Z}_x$ 是代数叠；见《代数叠》引理
\ref{algebraic-lemma-equivalent}。
由于 $\mathcal{Z}_x$ 通过该等价继承 $\mathcal{Z}$ 的所有性质，
由引理 \ref{stacks-properties-lemma-residual-gerbe} 中的唯一性显然可知，
$\mathcal{Z}_x$ 是 $\mathcal{X}$ 在 $x$ 处的剩余胚。
\end{proof}

\begin{lemma}
\label{stacks-properties-lemma-residual-gerbe-functorial}
设 $f : \mathcal{X} \to \mathcal{Y}$ 是代数叠态射。
设 $x \in |\mathcal{X}|$ 的像为 $y \in |\mathcal{Y}|$。
若 $x$ 和 $y$ 的剩余胚
$\mathcal{Z}_x \subset \mathcal{X}$ 与
$\mathcal{Z}_y \subset \mathcal{Y}$ 均存在，
则 $f$ 诱导交换图表
$$
\xymatrix{
\mathcal{X} \ar[d]_f & \mathcal{Z}_x \ar[l] \ar[d] \\
\mathcal{Y} & \mathcal{Z}_y \ar[l]
}
$$
\end{lemma}

\begin{proof}
选取域 $k$ 和满、平坦、局部有限表示的态射
$\Spec(k) \to \mathcal{Z}_x$。由引理
\ref{stacks-properties-lemma-residual-gerbe-points}，态射
$\Spec(k) \to \mathcal{Y}$ 经 $\mathcal{Z}_y$ 分解。
因此 $\mathcal{Z}_x \times_\mathcal{Y} \mathcal{Z}_y$
是 $\mathcal{Z}_x$ 的非空子叠，故由引理
\ref{stacks-properties-lemma-monomorphism-into-point}，它等于 $\mathcal{Z}_x$。
\end{proof}

\begin{lemma}
\label{stacks-properties-lemma-residual-gerbe-isomorphic}
设 $f : \mathcal{X} \to \mathcal{Y}$ 是代数叠态射。
设 $x \in |\mathcal{X}|$ 的像为 $y \in |\mathcal{Y}|$。
假设 $x$ 和 $y$ 的剩余胚
$\mathcal{Z}_x \subset \mathcal{X}$ 与
$\mathcal{Z}_y \subset \mathcal{Y}$ 均存在，
并且存在 $x$ 的等价类中的态射
$\Spec(k) \to \mathcal{X}$，使得
$$
\Spec(k) \times_\mathcal{X} \Spec(k)
\longrightarrow
\Spec(k) \times_\mathcal{Y} \Spec(k)
$$
是同构。则 $\mathcal{Z}_x \to \mathcal{Z}_y$ 是同构。
\end{lemma}

\begin{proof}
设 $k'/k$ 是域扩张。则
$$
\Spec(k') \times_\mathcal{X} \Spec(k')
\longrightarrow
\Spec(k') \times_\mathcal{Y} \Spec(k')
$$
是引理中的态射沿忠实平坦态射
$\Spec(k' \otimes k') \to \Spec(k \otimes k)$ 的基变换。
因此，引理所述性质与 $x$ 的等价类中态射
$\Spec(k) \to \mathcal{X}$ 的选择无关。
故可以假设 $\Spec(k) \to \mathcal{Z}_x$
满、平坦且局部有限表示。此时有
$$
\mathcal{Z}_x = [\Spec(k)/R]
$$
，其中 $R = \Spec(k) \times_\mathcal{X} \Spec(k)$；
见引理 \ref{stacks-properties-lemma-improve-unique-point} 的证明。
由于还有 $R = \Spec(k) \times_\mathcal{Y} \Spec(k)$，
可知引理 \ref{stacks-properties-lemma-residual-gerbe-functorial} 中的态射
$\mathcal{Z}_x \to \mathcal{Z}_y$ 全忠实；这是由
《代数叠》引理
\ref{algebraic-lemma-map-space-into-stack} 得到的。
例如，再由引理 \ref{stacks-properties-lemma-residual-gerbe-unique} 即得结论。
\end{proof}








\section{叠的维数}
\label{stacks-properties-section-dimension}

\noindent
利用代数空间在一点处的维数概念（《空间的性质》定义
\ref{spaces-properties-definition-dimension-at-point}），
可以定义代数叠 $\mathcal{X}$ 在点 $x$ 处的维数。
下一个引理中的结果可能是 $\infty$：这或者因为
$\mathcal{X}$ 不拟紧，或者因为遇到《例》第
\ref{examples-section-Noetherian-infinite-dimension} 节所述现象。

\begin{lemma}
\label{stacks-properties-lemma-dimension-at-point-well-defined}
设 $\mathcal{X}$ 是概形 $S$ 上局部 Noether 的代数叠，
且 $x \in |\mathcal{X}|$ 是 $\mathcal{X}$ 的点。
设 $[U/R] \to \mathcal{X}$ 是一个表示
（《代数叠》定义 \ref{algebraic-definition-presentation}），
其中 $U$ 是概形。设 $u \in U$ 是映到 $x$ 的点。
令 $e : U \to R$ 为“恒等”映射，并令 $s : R \to U$
为“源”映射；后者是代数空间的光滑态射。
令 $R_u$ 为 $s : R \to U$ 在 $u$ 上的纤维。则元素
$$
\dim_x(\mathcal{X}) = \dim_u(U) - \dim_{e(u)}(R_u) \in \mathbf{Z} \cup \infty
$$
与表示和 $x$ 上方之点 $u$ 的选择无关。
\end{lemma}

\begin{proof}
由于 $R \to U$ 光滑，概形 $R_u$ 在 $\kappa(u)$ 上光滑，
因而具有有限维数。另一方面，概形 $U$ 局部 Noether，
但这并不保证 $\dim_u(U)$ 有限。因此，该差属于
$\mathbf{Z} \cup \{\infty\}$。

\medskip\noindent
设 $[U'/R'] \to \mathcal{X}$ 和 $u' \in U'$ 为第二组表示与点，
其中 $U'$ 是概形，且 $u'$ 映到 $x$。
考虑代数空间 $P = U \times_\mathcal{X} U'$。
由引理 \ref{stacks-properties-lemma-points-cartesian}，存在
$p \in |P|$ 映到 $u$ 和 $u'$。由于 $P \to U$ 和
$P \to U'$ 光滑，可得
$\dim_p(P) = \dim_u(U) + \dim_p(P_u)$ 以及
$\dim_p(P) = \dim_{u'}(U') + \dim_p(P_{u'})$；
见《空间的态射》引理
\ref{spaces-morphisms-lemma-smoothness-dimension-spaces}。
注意
$$
R'_{u'} = \Spec(\kappa(u')) \times_\mathcal{X} U'
\quad\text{且}\quad
P_u = \Spec(\kappa(u)) \times_\mathcal{X} U'
$$
用态射 $\Spec(\Omega) \to P$ 表示 $p \in |P|$。
由于 $p$ 同时映到 $u$ 和 $u'$，它诱导下列两个复合之间的 $2$-态射：
$\Spec(\Omega) \to \Spec(\kappa(u')) \to \mathcal{X}$ 和
$\Spec(\Omega) \to \Spec(\kappa(u)) \to \mathcal{X}$
，进而定义同构
$$
\Spec(\Omega) \times_{\Spec(\kappa(u'))} R'_{u'}
\cong
\Spec(\Omega) \times_{\Spec(\kappa(u))} P_u
$$
，它是 $\Spec(\Omega)$ 上代数空间的同构，并将
$\Omega$-有理点 $(1, e'(u'))$ 映到 $(1, p)$（略去一些细节）。
由此可得
$$
\dim_{e'(u')}(R'_{u'}) = \dim_p(P_u)
$$
，这是由《空间的态射》引理
\ref{spaces-morphisms-lemma-dimension-fibre-after-base-change}.
得到的。由对称性，有
$\dim_{e(u)}(R_u) = \dim_p(P_{u'})$。
综合以上各式，即得与选择无关。
\end{proof}

\noindent
可以利用上述引理作如下定义。

\begin{definition}
\label{stacks-properties-definition-dimension-at-point}
设 $\mathcal{X}$ 是概形 $S$ 上局部 Noether 的代数叠，
且 $x \in |\mathcal{X}|$ 是 $\mathcal{X}$ 的点。
设 $[U/R] \to \mathcal{X}$ 是一个表示
（《代数叠》定义 \ref{algebraic-definition-presentation}），
其中 $U$ 是概形；并设 $u \in U$ 是映到 $x$ 的点。
将{\it $\mathcal{X}$ 在 $x$ 处的维数}定义为元素
$\dim_x(\mathcal{X}) \in \mathbf{Z} \cup \infty$，满足
$$
\dim_x(\mathcal{X}) = \dim_u(U)-\dim_{e(u)}(R_u).
$$
记号同引理 \ref{stacks-properties-lemma-dimension-at-point-well-defined}。
\end{definition}

\noindent
当 $\mathcal{X}$ 是概形时（《拓扑》定义
\ref{topology-definition-Krull}），叠在一点处的维数与通常概念一致；
更一般地，当 $\mathcal{X}$ 是局部 Noether 代数空间时亦如此
（《空间的性质》定义
\ref{spaces-properties-definition-dimension-at-point}）。

\begin{definition}
\label{stacks-properties-definition-dimension}
设 $S$ 是概形，且 $\mathcal{X}$ 是 $S$ 上局部 Noether 的代数叠。
将 $\mathcal{X}$ 的{\it 维数} $\dim(\mathcal{X})$ 定义为
$$
\dim(\mathcal{X}) = \sup\nolimits_{x \in |\mathcal{X}|} \dim_x(\mathcal{X})
$$
\end{definition}

\noindent
若 $\mathcal{X}$ 是概形
（《性质》引理 \ref{properties-lemma-dimension}）
或代数空间（《空间的性质》定义
\ref{spaces-properties-definition-dimension}），
该维数定义与通常概念一致。

\begin{remark}
\label{stacks-properties-remark-dimension-empty-stack}
若 $\mathcal{X}$ 是域上非空的有限型叠，则
$\dim(\mathcal{X})$ 是整数。对任意局部 Noether 代数叠
$\mathcal{X}$，$\dim(\mathcal{X})$ 属于
$Z\cup \{\pm \infty\}$；
且 $\dim(\mathcal{X}) = -\infty$ 当且仅当 $\mathcal{X}$ 为空。
\end{remark}

\begin{example}
\label{stacks-properties-example-dimension-quotient-stack}
设 $X$ 是域 $k$ 上的有限型概形，且 $G$ 是作用在 $X$ 上的
$k$ 上有限型群概形。则商叠 $[X/G]$ 的维数等于
$\dim(X)-\dim(G)$。特别地，分类叠
$BG=[\Spec(k)/G]$ 的维数为 $-\dim(G)$。
因此，与概形或代数空间的情形不同，代数叠的维数可以是负整数。
\end{example}






\section{局部不可约性}
\label{stacks-properties-section-irreducible-local-ring}

\noindent
我们已经在《性质》第 \ref{properties-section-unibranch} 节
定义了概形在一点处的几何分支数，并在《空间的性质》第
\ref{spaces-properties-section-irreducible-local-ring} 节
定义了代数空间在一点处的几何分支数。
设 $n \in \mathbf{N}$。对局部环 $A$，令
$$
P_n(A) = \text{the number of geometric branches of }A\text{ is }n
$$
对光滑环同态 $A \to B$ 和 $B$ 中位于 $A$ 的素理想
$\mathfrak p$ 上方的素理想 $\mathfrak q$，有
$$
P_n(A_\mathfrak p) \Leftrightarrow P_n(B_\mathfrak q)
$$
，这是由《代数进阶》引理
\ref{more-algebra-lemma-invariance-number-branches-smooth} 得到的。
如同《空间的性质》注
\ref{spaces-properties-remark-list-properties-local-ring-local-etale-topology}，
可以利用 $P_n$ 定义概形芽 $(U, u)$ 的 \'etale 局部性质
$\mathcal{P}_n$，令
$\mathcal{P}_n(U, u) = P_n(\mathcal{O}_{U, u})$。
代数空间 $X$ 在点 $x$ 处的相应性质 $\mathcal{P}_n$
（见《空间的性质》定义
\ref{spaces-properties-definition-property-at-point}）
恰好是性质“$X$ 在 $x$ 处的几何分支数为 $n$”；
见《空间的性质》定义
\ref{spaces-properties-definition-number-of-geometric-branches}。
此外，性质 $\mathcal{P}_n$ 是光滑局部的；
见《下降》定义 \ref{descent-definition-local-at-point}。
这或者由上面显示的等价式推出，或者由《态射进阶》引理
\ref{more-morphisms-lemma-number-of-branches-and-smooth} 推出。
因此，定义 \ref{stacks-properties-definition-property-at-point} 适用，
从而得到代数叠在一点处的相应概念。

\begin{definition}
\label{stacks-properties-definition-number-of-geometric-branches}
设 $\mathcal{X}$ 是代数叠，且 $x \in |\mathcal{X}|$。
\begin{enumerate}
\item 若引理 \ref{stacks-properties-lemma-local-source-target-at-point} 的等价条件
对上述定义的 $\mathcal{P}_n$ 成立，则
{\it $\mathcal{X}$ 在 $x$ 处的几何分支数}为
$n \in \mathbf{N}$；否则为 $\infty$。
\item 若 $\mathcal{X}$ 在 $x$ 处的几何分支数为 $1$，
则称 $\mathcal{X}$ {\it 在 $x$ 处几何单分支}。
\end{enumerate}
\end{definition}






\section{有限性条件与点}
\label{stacks-properties-section-points-conditions}

\noindent
本节是《适度空间》第
\ref{decent-spaces-section-points-monomorphisms} 节
对代数叠之点的类似版本。

\begin{lemma}
\label{stacks-properties-lemma-UR-quasi-compact-above-x}
设 $\mathcal{X}$ 是代数叠，且 $x \in |\mathcal{X}|$ 是一个点。
下列条件等价：
\begin{enumerate}
\item $x$ 的等价类中的某个态射
$\Spec(k) \to \mathcal{X}$ 拟紧；
\item $x$ 的等价类中的任一态射
$\Spec(k) \to \mathcal{X}$ 拟紧。
\end{enumerate}
\end{lemma}

\begin{proof}
设 $\Spec(k) \to \mathcal{X}$ 属于 $x$ 的等价类，
并设 $k'/k$ 是域扩张。需要证明
$\Spec(k) \to \mathcal{X}$ 拟紧，当且仅当
$\Spec(k') \to \mathcal{X}$ 拟紧。
这由《空间的态射》引理
\ref{spaces-morphisms-lemma-surjection-from-quasi-compact}
以及《代数叠》引理
\ref{algebraic-lemma-representable-transformations-property-implication}
的原理推出。
\end{proof}

\noindent
有时将满足引理 \ref{stacks-properties-lemma-UR-quasi-compact-above-x}
之等价条件的点 $x \in |\mathcal{X}|$
称为“{\it 拟紧点}”。
